\documentclass[10pt,twocolumn,letterpaper]{article}

\usepackage[review]{cvpr}
\usepackage{amssymb}
\captionsetup{font=footnotesize}
\definecolor{cvprblue}{rgb}{0.21,0.49,0.74}
\usepackage[pagebackref,breaklinks,colorlinks,allcolors=cvprblue]{hyperref}
\def\paperID{*****}
\def\confName{CVPR}
\def\confYear{2027}

% ============================================================================
% numbers.tex -- every number that appears in the paper, with its source.
% RULE: no numeral enters main.tex directly. It is defined here, and the
% comment above it names the experiment directory that produced it. A number
% with no experiment behind it is written as \TODOnum so it cannot be mistaken
% for a measurement.
% ============================================================================

\newcommand{\TODOnum}{\textcolor{red}{\textbf{[unmeasured]}}}

% --- E00: DSEC exposure survey (experiments/e00_exposure_survey) -------------
% 18 train sequences, 21142 frames, measured from the published
% <seq>_image_exposure_timestamps_left.txt files.
\newcommand{\dsecExpMin}{118}              % us, zurich_city_06_a
\newcommand{\dsecExpMax}{14\,996}          % us, the auto-exposure ceiling
\newcommand{\dsecExpRatio}{127}            % max/min, rounded
\newcommand{\dsecFramePeriod}{50}          % ms, every sequence
\newcommand{\dsecPinnedFrac}{32.5}         % percent of frames at the ceiling
\newcommand{\dsecPinnedCount}{6866}
\newcommand{\dsecFrameCount}{21\,142}
\newcommand{\dsecCeilingSeqs}{six}         % sequences pinned at the ceiling
\newcommand{\dsecOtherCeilingSeqs}{five}   % the ceiling sequences not measured here
\newcommand{\dsecWithinSeqMin}{337}        % us, interlaken_00_d
\newcommand{\dsecWithinSeqMax}{4207}       % us, interlaken_00_d
\newcommand{\dsecWithinSeqRatio}{12.5}
\newcommand{\dsecCotriggerFrames}{7080}    % frames checked, L/R start identical
\newcommand{\dsecCotriggerPct}{100}
\newcommand{\dsecMidDiffMax}{190}          % us, max |mid_L - mid_R|
\newcommand{\dsecWidthDiffMax}{380}        % us, max |w_L - w_R|
\newcommand{\dsecTsResidual}{1}            % us, |image_ts - avg(mid_L,mid_R)|

% --- E01: events inside one exposure (experiments/e01_intra_exposure_change) -
\newcommand{\dayExpUs}{1478}               % interlaken_00_c, median
\newcommand{\dayEvents}{20\,616}
\newcommand{\dayPixFrac}{6.0}              % percent
\newcommand{\dayCrossings}{1.10}
\newcommand{\dayFrames}{536}
\newcommand{\nightExpUs}{14\,996}          % zurich_city_09_a
\newcommand{\nightEvents}{310\,320}
\newcommand{\nightPixFrac}{38.2}           % percent
\newcommand{\nightCrossings}{2.60}
\newcommand{\nightFrames}{452}
\newcommand{\nightDayEventRatio}{15.1}

% --- E02: RVT window alignment (experiments/e02_rvt_window_alignment) --------
% Read from the released source, not from the paper's description of it.
\newcommand{\rvtWindow}{50}                % ms, stacked_histogram_dt=50
\newcommand{\rvtBins}{ten}                 % bins per window
\newcommand{\rvtBinWidth}{5}               % ms per bin
\newcommand{\rvtUniformCentroid}{-25}      % ms, uniform-weight reference point,
                                           % superseded as a claim by E17.

% --- E04: RVT loads under torch 2.7.1 (experiments/e04_rvt_port_risk) --------
\newcommand{\rvtParams}{4.41}              % M, rvt-t
\newcommand{\rvtMissingKeys}{0}
\newcommand{\rvtUnexpectedKeys}{0}

% --- E06: DSEC-Det label forensics (experiments/e06_dsec_det_label_forensics)-
\newcommand{\ddBoxes}{390\,118}
\newcommand{\ddTracks}{6693}
\newcommand{\ddSeqs}{60}
\newcommand{\ddSpacingMedian}{50\,001}     % us
\newcommand{\ddSubFrameFrac}{0.0000}       % fraction of spacings below 40 ms
\newcommand{\ddSubFrameFracPct}{0.00}      % the same, as a percentage
\newcommand{\ddInterpResidual}{0.707}      % px, median second-difference
\newcommand{\ddInterpResidualPninety}{1.819}
\newcommand{\ddInterpBelowThresh}{7.1}     % percent below 0.01 px
\newcommand{\ddTriples}{361\,628}

% --- E07: LET-3D-AP citation graph (experiments/e07_let3dap_citation_graph) --
\newcommand{\letCitations}{40}             % returned by the API, coverage unknown
\newcommand{\letTemporalHits}{6}
\newcommand{\letAxisHits}{0}               % none apply the decomposition to time

% --- Removed 2026-09-05: quantities awaiting detection runs this paper does not
% contain. No detection score is computed anywhere in the paper or the
% supplement, so the tolerant-metric and re-anchoring macros below have no text
% that could use them and are left undefined rather than rendering a placeholder.
% \newcommand{\tauHatCI}{\TODOnum}
% \newcommand{\sigmaTau}{\TODOnum}
% \newcommand{\panelCSlope}{\TODOnum}
% \newcommand{\panelCIntercept}{\TODOnum}
% \newcommand{\apSyncDelta}{\TODOnum}
% \newcommand{\apBiasDelta}{\TODOnum}

%% ===== MEASURED 2026-09-02, CORRECTED 2026-09-03 ============================
%% E09+E11 per-object event-centroid dispersion inside one exposure. (E12 was
%% retracted on 2026-09-05, every value in this block is E11's.)
%% No detector, no checkpoint, no velocity denominator, no assumed center noise.
%% The values below are E11's exact-window re-run (result_exact.json), which
%% restricts events to [a,b] in microseconds. The first run sliced to
%% whole-millisecond boundaries and is superseded, it produced 208.8 us.
\newcommand{\sigmaEvtNight}{206.6}   %% us within-frame sd, night, all pixels
\newcommand{\sigmaEvtNullNight}{93.0}   %% us analytic uniform null
\newcommand{\sigmaEvtExcessNight}{183.6}   %% us excess over the null
\newcommand{\sigmaEvtFramesNight}{1134}   %% frames measured
\newcommand{\sigmaEvtFracNight}{0.9012}   %% fraction of frames sd>null
%% E11 flicker controls, exact window (result_exact.json)
\newcommand{\excessClean}{180.6}   %% us phase-locked pixels excluded
\newcommand{\excessRandCtl}{182.9}   %% us matched random exclusion
\newcommand{\lockedSensorPct}{4.67}   %% percent OF THE SENSOR, not of live pixels
%% E12 RETRACTED 2026-09-03 and removed from the paper 2026-09-04. An integer
%% window cancels a periodic source's zeroth moment, not its first: the centroid
%% contribution is -(A/lambda_0) sin(phi)/omega. The widths below were measured
%% correctly, the cancellation argument they were reported under was false, so
%% the macros are left undefined and no width appears in the paper.
%% \newcommand{\excessHalfPeriod}{212.6}   %% us w=5000 (0.5 periods)
%% \newcommand{\excessOnePeriod}{195.2}   %% us w=10000 (1.0 periods)
%% \newcommand{\excessFullExp}{183.6}   %% us w=14996 (1.4996 periods)
%% \newcommand{\excessTwoPeriod}{246.8}   %% us w=20000 (2.0 periods)
%% E08 label noise whiteness
\newcommand{\labelAcfMeasured}{-0.3566}   %% D3 acf lag1 measured on real labels
\newcommand{\labelAcfWhite}{-0.740}   %% same code on white noise
\newcommand{\labelAcfTheory}{-0.75}   %% analytic white-noise value
\newcommand{\ntracksNoise}{5250}
\newcommand{\nthirdDiff}{629\,442}
%% E10 mains flicker
\newcommand{\flickerLineHz}{100}   %% Hz, strongest night line
\newcommand{\flickerZmedian}{2.935}   %% night 100Hz Rayleigh Z, null 0.693
\newcommand{\flickerZnull}{0.693}
\newcommand{\flickerFracLocked}{0.2407}   %% fraction of LIVE pixels p<1e-3
\newcommand{\flickerCtrlNightZ}{0.581}   %% night 137Hz off-frequency control
\newcommand{\flickerCtrlDayZ}{0.719}   %% day 100Hz control

%% --- E00, shape of the exposure distribution (21142 train frames) ----------
\newcommand{\dsecExpGapLo}{4207}           % us, largest non-ceiling exposure
\newcommand{\dsecExpGapHi}{14\,996}        % us, the ceiling
\newcommand{\dsecExpGapCount}{zero}        % frames strictly between the two
\newcommand{\pctSixtySix}{66}              % the percentile index itself
\newcommand{\pctSixtyEight}{68}
\newcommand{\pctBandLo}{5}                 % Fig. 1's whisker band
\newcommand{\pctBandHi}{95}
\newcommand{\dsecExpPsixtysix}{2999}       % us, percentile 66
\newcommand{\dsecExpPsixtyeight}{14\,996}  % us, percentile 68

%% --- E06, the per-box time field of the released labels --------------------
\newcommand{\ddDistinctT}{70\,379}         % distinct values of t, summed within sequence
\newcommand{\ddBoxesPerT}{5.54}            % boxes per distinct value

%% --- E09, the measurement and its strata -----------------------------------
\newcommand{\sigmaEvtObjNight}{6}          % median qualifying objects per frame
\newcommand{\sigmaEvtSdPninetyNight}{344.5}% us, p90 of the within-frame sd (exact window)
\newcommand{\sigmaEvtPtpNight}{563.9}      % us, median within-frame peak to peak
\newcommand{\sigmaEvtDay}{18.0}            % us, interlaken_00_c within-frame sd
\newcommand{\sigmaEvtNullDay}{21.6}        % us, its analytic null
\newcommand{\sigmaEvtFramesDay}{20}        % usable frames, day
\newcommand{\sigmaEvtFracDay}{0.35}        % fraction of day frames sd>null
\newcommand{\excessDay}{0}                 % us, the day excess: sd is below the null
\newcommand{\sigmaEvtExpDay}{1556}         % us, median exposure of the day frames used
\newcommand{\sdRandBoxes}{338.2}           % us, random-position boxes, same shapes
\newcommand{\nullRandBoxes}{134.5}         % us, their analytic null
\newcommand{\excessRandBoxes}{306.3}       % us, their excess -- larger than the labels'
\newcommand{\framesRandBoxes}{1123}        % zurich_city_09_a.json, n_frames_with_rand
\newcommand{\minEvents}{200}               % events per box required
\newcommand{\minObjects}{3}                % qualifying objects per frame required

%% --- E09, the per-track structure of the frame-centered residual ------------
%% experiments/e09_per_object_evidence_time/per_track.py, re-run 2026-09-03 on
%% the corrected exposure window.  Computed after removing each frame's own
%% mean, so the frame-level component that a declared scalar timestamp would
%% cancel is removed by construction.
\newcommand{\evtTracks}{103}               % tracks with at least five observations
\newcommand{\betweenTrackSd}{149.2}        % us, between-track sd of the residual
\newcommand{\betweenTrackNull}{39.3}       % us, expected with no track effect
\newcommand{\trackVarRatio}{14.4}          % observed / no-effect variance ratio
\newcommand{\trackVarRatioCell}{8.6}       % same ratio after removing 32x32 cell means
\newcommand{\trackAcf}{0.721}              % lag-1 autocorrelation along a track
\newcommand{\acfLowDisp}{0.611}            % tracks below the displacement median
\newcommand{\acfHighDisp}{0.773}           % tracks above it
\newcommand{\dispSplit}{278}               % px, the median path length

%% --- E10, the mains line ---------------------------------------------------
\newcommand{\flickerPeriodUs}{10\,000}     % us, one 100 Hz period
\newcommand{\flickerPeakRatio}{10\,841}    % night peak over the 60-160 Hz continuum
\newcommand{\flickerLivePixNight}{35\,990} % live night pixels tested
\newcommand{\flickerLivePixDay}{17\,800}   % live day pixels tested
\newcommand{\flickerZpNinetyfive}{16.52}   % night 100 Hz p95 of Z
\newcommand{\flickerNullPninetyfive}{3.00} % Exp(1) p95
\newcommand{\flickerNullRate}{0.001}       % Exp(1) rate at p<1e-3
\newcommand{\flickerFracLockedSix}{0.0750} % fraction of live night pixels, p<1e-6
\newcommand{\flickerFracLockedDay}{0.0008} % day 100 Hz, same test

%% --- E11, the three pixel sets (result_exact.json) -------------------------
\newcommand{\excessMaskPix}{14\,334}       % phase-locked pixels excluded
\newcommand{\sdAll}{206.6}
\newcommand{\nullAll}{93.0}
\newcommand{\fracAll}{0.9012}
\newcommand{\sdClean}{205.1}
\newcommand{\nullClean}{97.1}
\newcommand{\fracClean}{0.8942}
\newcommand{\sdRandCtl}{206.3}
\newcommand{\nullRandCtl}{94.8}
\newcommand{\fracRandCtl}{0.8967}
\newcommand{\framesEleven}{1134}           % ALL and CLEAN arms
\newcommand{\framesRandCtl}{1133}          % matched random exclusion

%% --- E12, four analysis-window widths, RETRACTED with the block above -------
%% \newcommand{\wHalfPeriod}{5000}      %% \sdHalfPeriod 218.1  \nullHalfPeriod 46.5
%% \newcommand{\wOnePeriod}{10\,000}    %% \sdOnePeriod  211.5  \nullOnePeriod  74.4
%% \newcommand{\wFullExp}{14\,996}      %% \sdFullExp    206.6  \nullFullExp    93.0
%% \newcommand{\wTwoPeriod}{20\,000}    %% \sdTwoPeriod  271.2  \nullTwoPeriod 108.0
%% frames 1085 / 1118 / 1134 / 1137, frac over null 0.988 / 0.961 / 0.901 / 0.945

%% --- E08, object speed on DSEC-Det -----------------------------------------
\newcommand{\labelSpeedMedian}{40}         % px/s, median box-center speed
\newcommand{\labelSpeedPninety}{150}       % px/s, p90
\newcommand{\labelSpeedPninetynine}{360}   % px/s, p99
\newcommand{\offsetAtMedian}{1.0}          % px, 25 ms at the median speed
\newcommand{\offsetAtPninetynine}{9.0}     % px, 25 ms at p99

%% --- E13, the label center noise, solved rather than bounded ---------------
%% Difference-order ladder plus a two-rate elimination, the ladder is still
%% decreasing at k=7, so the total is a geometric extrapolation.
\newcommand{\sigmaCtotal}{0.4974}          % px, total center noise
\newcommand{\sigmaCquant}{0.1443}          % px, the integer-lattice component
\newcommand{\sigmaCannot}{0.476}           % px, the annotation component
\newcommand{\sigmaClattice}{0.5}           % px, the lattice step DSEC-Det stores on
\newcommand{\sigmaCmillis}{11.9}           % ms, sigmaCannot at the median speed
\newcommand{\ladderK}{7}                   % highest difference order taken

%% --- E14, the dispersion on the pixel scale --------------------------------
%% One multiplication: displacement = dispersion x object speed.
\newcommand{\pxAtMedian}{0.0073}           % px, at the median object speed
\newcommand{\pxAtPninety}{0.0275}          % px, at p90
\newcommand{\pxAtPninetynine}{0.0661}      % px, at p99
\newcommand{\pxFracMedian}{1.5}            % percent of sigmaCtotal
\newcommand{\pxFracPninetynine}{13.3}      % percent of sigmaCtotal
\newcommand{\speedToQuant}{786}            % px/s, speed at which it reaches sigmaCquant
\newcommand{\speedToNoise}{2709}           % px/s, speed at which it reaches sigmaCtotal

%% --- E16, the common offset from mid-exposure ------------------------------
\newcommand{\offsetAllEvents}{\ensuremath{-16.8}}    % us, night, all events
\newcommand{\offsetInBoxes}{\ensuremath{-102.3}}     % us, night, inside labeled boxes
\newcommand{\offsetDayBound}{2}            % us, both day figures fall below this
\newcommand{\offsetFramesSixteen}{1811}    % night frames surveyed
\newcommand{\pxOffsetPninetynine}{0.037}   % px, offsetInBoxes at p99 speed

%% --- macros for quantities that were still inline numerals ------------------
\newcommand{\dsecSeqCount}{eighteen}       % E00, sequences whose exposure file resolves
\newcommand{\dsecFrameRate}{20}            % Hz, label and frame rate
\newcommand{\ddArchiveBytes}{4\,655\,023}  % bytes, dsec-det_left_object_detections.zip
\newcommand{\ddArchiveDate}{2026-09-01}    % retrieval date
\newcommand{\sigmaEvtFramesLabelled}{1557} % E09, labelled frames in the night sequence
\newcommand{\mainsHz}{50}                  % Hz, Swiss mains
\newcommand{\flickerBandLo}{90}            % Hz, search band for the strongest line
\newcommand{\flickerBandHi}{110}
\newcommand{\flickerContLo}{60}            % Hz, continuum band
\newcommand{\flickerContHi}{160}
\newcommand{\flickerPeakRatioDay}{6.6}     % day peak over the same continuum
\newcommand{\flickerMinEvents}{100}        % events per pixel required
\newcommand{\flickerCtrlNightRate}{0.0004} % off-frequency arm, rate at p<1e-3
\newcommand{\periodsFullExp}{1.4996}   % E10, mains periods inside the ceiling exposure
\newcommand{\rvtCentroidMag}{25}           % ms, |rvtUniformCentroid|
\newcommand{\rvtTwoBranchScale}{7.2}       % ms, rvtCentroidMag/sqrt(12)

%% ===== E17 measured 2026-09-03: the network's temporal influence profile ====
%% 480 samples of real gen1 val, released rvt-t, recurrent state warmed eight
%% steps, each of the ten bins occluded in turn. This converts E02's
%% uniform-weight REFERENCE POINT into a MEASUREMENT.
\newcommand{\rvtCentroidMeas}{\ensuremath{-23.81}}   %% ms, influence-weighted centroid,
                                           %% E45. E17 and E34 reported -24.94 because the
                                           %% occluded pass received the recurrent state the
                                           %% reference pass RETURNED, not the one entering
                                           %% the step, E45 repeats them with the entering
                                           %% state and supersedes both.
\newcommand{\rvtCentroidMagMeas}{23.81}    %% ms, its magnitude, for "before"
\newcommand{\rvtCentroidUnif}{\ensuremath{-25.00}}   %% ms, uniform-weight reference
\newcommand{\rvtCentroidDiff}{\ensuremath{+1.19}}    %% ms, measured minus uniform
\newcommand{\rvtCentroidDiffPct}{23.8}     %% percent of one 5 ms bin: 1.190/5.0
\newcommand{\rvtInfluenceSamples}{480}
\newcommand{\rvtWarmSteps}{eight}          %% recurrent warm-up steps
\newcommand{\rvtInfluenceCV}{0.131}   %% coefficient of variation across the ten bins
\newcommand{\rvtInfluenceMaxMin}{1.52} %% heaviest bin over lightest
\newcommand{\rvtHalfOlder}{48.6}       %% percent of influence in the five older bins
\newcommand{\rvtHalfNewer}{51.4}       %% percent in the five newer bins
\newcommand{\tauHatRVT}{\rvtCentroidMeas}  %% E17, no longer a placeholder

%% --- proportions, written as percentages everywhere in the paper -----------
%% Same measurements as the fractions above, one format is used throughout so
%% that adjacent proportions are comparable by eye.
\newcommand{\sigmaEvtFracNightPct}{90.1}   % E11 exact, ALL arm
\newcommand{\sigmaEvtFracDayPct}{35}       % E09 day arm
\newcommand{\fracAllPct}{90.1}
\newcommand{\fracCleanPct}{89.4}
\newcommand{\fracRandCtlPct}{89.7}
\newcommand{\flickerFracLockedPct}{24.07}      % E10, of LIVE pixels
\newcommand{\flickerFracLockedDayPct}{0.08}    % E10, day arm
\newcommand{\flickerFracLockedSixPct}{7.50}    % E10, p<1e-6
\newcommand{\flickerCtrlNightRatePct}{0.04}    % E10, off-frequency arm
\newcommand{\flickerNullRatePct}{0.1}          % Exp(1) rate at p<1e-3
\newcommand{\cellGrid}{32\,$\times$\,32}       % E09, image-cell granularity

%% --- the ratio of the two temporal terms (E17 over E11) -------------------
\newcommand{\termRatio}{130}               % 23.81 ms / 183.6 us
\newcommand{\pxPredMedian}{1.0}            % px, tauHatRVT at the median speed
\newcommand{\pxPredPninetynine}{9.0}       % px, tauHatRVT at p99

%% ===== E19 measured 2026-09-03: the scalar dispersion against each box's own
%% mains modulation (experiments/e19_modulation_strata). For each box the
%% modulation depth at 100 Hz is m = 2|C| with C = (1/N) sum exp(-i w t_k) over
%% its own events, the within-frame dispersion is then recomputed inside
%% quartiles of m. This is what demotes the scalar dispersion to a bounded
%% statement: it rises monotonically with m and no unmodulated stratum exists.
\newcommand{\modFrames}{1134}              %% E19, frames entering the strata
\newcommand{\modBoxes}{6841}               %% E19, boxes entering the strata
\newcommand{\modPten}{0.355}               %% E19, p10 of m at 100 Hz
\newcommand{\modMedian}{0.467}             %% E19, median of m at 100 Hz
\newcommand{\modPninety}{0.609}            %% E19, p90 of m at 100 Hz
\newcommand{\modFloor}{0.125}              %% E19, m for 200 events at random phase
\newcommand{\excessModQone}{93.1}          %% us, E19, lowest 100 Hz quartile
\newcommand{\excessModQtwo}{128.1}         %% us, E19, second quartile
\newcommand{\excessModQthree}{131.5}       %% us, E19, third quartile
\newcommand{\excessModQfour}{208.0}        %% us, E19, highest quartile
\newcommand{\sdModQone}{119.3}             %% us, E19, sd of the lowest quartile
\newcommand{\sdModQtwo}{149.1}
\newcommand{\sdModQthree}{156.6}
\newcommand{\sdModQfour}{227.2}
\newcommand{\nullModQone}{74.6}            %% us, E19, its analytic null
\newcommand{\nullModQtwo}{76.2}
\newcommand{\nullModQthree}{85.1}
\newcommand{\nullModQfour}{91.5}
\newcommand{\framesModQone}{270}           %% E19, frames per stratum
\newcommand{\framesModQtwo}{243}
\newcommand{\framesModQthree}{207}
\newcommand{\framesModQfour}{265}
\newcommand{\modQoneLo}{0.122}             %% E19, the quartile edges of m
\newcommand{\modQoneHi}{0.410}
\newcommand{\modQtwoHi}{0.467}
\newcommand{\modQthreeHi}{0.534}
\newcommand{\modQfourHi}{1.050}
\newcommand{\modExcessRatio}{2.2}          %% E19, Q4 over Q1 at 100 Hz
\newcommand{\excessShamQone}{150.3}        %% us, E19, 137 Hz control, lowest quartile
\newcommand{\excessShamQfour}{187.6}       %% us, E19, 137 Hz control, highest quartile
\newcommand{\shamExcessRatio}{1.25}        %% E19, its Q4 over Q1
\newcommand{\flickerCtrlHz}{137}           %% Hz, the off-frequency of E10 and E19

%% ===== E20 RETRACTED 2026-09-05 ==============================================
%% The spatial gradient of evidence time. Five controlled passes, E25 generalized
%% the measurement to all six ceiling sequences and found the alignment STRONGER
%% where objects move slower, which a per-object motion gradient cannot produce,
%% and E26 compared each box against an equal-area annulus of its own surrounding
%% background: in four of six sequences the background aligns better than the
%% object, and eroding the box to its object-dominated core makes it worse in five
%% of six. The signal is ego-motion. Every macro of this block is left undefined so
%% that no sentence of the paper can use one. See experiments/e20_spatial_gradient
%% and docs/PROTOCOL_LEDGER.md entry 6.
%% \newcommand{\gradBoxes}{6627}              %% E20, qualifying boxes
%% \newcommand{\gradMinEvents}{400}           %% E20, events per box required
%% \newcommand{\gradSpeedFloor}{5}            %% px/s, E20, label-speed floor
%% \newcommand{\gradSpeedMedian}{58.3}        %% px/s, E20, median label speed
%% \newcommand{\gradCos}{\ensuremath{+0.152}} %% E20, mean cos(grad t, v)
%% \newcommand{\gradCosMedian}{\ensuremath{+0.304}}   %% E20, its median
%% \newcommand{\gradCosNullVel}{\ensuremath{+0.028}}  %% E20, shuffled-velocity null
%% \newcommand{\gradCosNullTime}{\ensuremath{+0.020}} %% E20, time-shuffled null
%% \newcommand{\gradSE}{0.012}                %% E20, 1/sqrt(n), a bound on the SE
%% \newcommand{\gradSigmas}{12.4}             %% E20, measured value in those units
%% \newcommand{\gradNullSigmas}{two}          %% E20, both nulls lie within this many
%% \newcommand{\gradMag}{11.23}               %% us/px, E20, median |grad t|
%% \newcommand{\gradMagShuffle}{5.47}         %% us/px, E20, its time-shuffled floor
%% \newcommand{\gradMagRatio}{2.05}           %% E20, measured over floor
%% \newcommand{\gradTranslatePx}{0.87}        %% px, gradSpeedMedian over dsecExpMax
%% \newcommand{\gradBoxSide}{42}              %% px, median sqrt(wh) of the released

%% --- the frame drawn in the qualitative figure -----------------------------
%% src/make_fig_qualitative.py, which ranks the sequence's candidate frames by
%% their within-frame dispersion and draws the MEDIAN one.
\newcommand{\figFiveCands}{602}            %% candidate frames ranked
\newcommand{\figFiveBoxes}{six}            %% qualifying boxes in the drawn frame
\newcommand{\figFiveEvents}{407\,942}      %% events inside its exposure window
\newcommand{\figFiveSd}{220.3}             %% us, its within-frame sd

%% ===== RETRACTED 2026-09-07: the E21, E23 and E24 blocks.
%% E24 fitted |e_par| and |e_perp|, E[|eps+tau v|] is not E[|eps|]+tau v, so its
%% coefficient difference does not identify a temporal offset (external review #3,
%% Critical 1, and the internal checklist audit of the same day). E23's acceleration
%% regressor was |dv|/dt, a magnitude entered where a signed a_par is required, and it
%% correlates with speed at +0.36. E21 and E24 both took the velocity from a FORWARD
%% difference, which shares the label-center noise with the residual and manufactures
%% a lag of tens of milliseconds when none is present (E29). All three are replaced by
%% E27/E30/E31/E32 below. The macros are removed rather than commented out so that no
%% text can reach a withdrawn number, see docs/REVIEW3_RESPONSE.md and
%% docs/PROTOCOL_LEDGER.md case 7.

%% ===== E27/E30/E31/E32 measured 2026-09-07: the effective OUTPUT time of the
%% released checkpoint, re-measured after three defects were found in E21 and E24.
%% Raw rows: experiments/e27_rows/rows.npz (35 523 matches, 406 sequences), one
%% GPU pass, every fit below is a CPU re-analysis of that one table.
%%   fit.log            the specification ladder and the sign-flip control
%%   placebo_diag.log   what the rotated-velocity coefficient is
%%   injection.log      recovery of an injected lag, and the hypothesis tests
%% Convention, stated once and used everywhere from here: p - g = -tau v, so a
%% POSITIVE tau means the output describes an EARLIER state than the label instant.
\newcommand{\outMatches}{22\,534}          %% E27, matched boxes with a centered velocity
\newcommand{\outRowsTotal}{35\,523}        %% E27, matches before requiring both neighbors
\newcommand{\outSeqs}{376}                 %% E27, sequences contributing to the fit
\newcommand{\outSeqsScanned}{406}          %% E27, sequences with at least one match
\newcommand{\outIoUfloor}{0.3}             %% E27, overlap required to match
\newcommand{\outIoUmedian}{0.873}          %% E27, median overlap of the matches
\newcommand{\outSpeedMedian}{5.4}          %% px/s, E27, median label speed, no floor
\newcommand{\outLinkPct}{79.0}             %% percent of Gen1 val boxes with a successor
                                           %% at IoU>=0.3, E36 (e27_rows/split_stats.log).
                                           %% The 81.5 here before was E21's floored subset.
\newcommand{\outLinkIoU}{0.3}              %% the overlap at which they link
\newcommand{\outSpeedSizeCorr}{\ensuremath{+0.139}} %% corr(|v|,sqrt(A)) on the E27 fit
                                           %% sample, E36. E21's +0.349 came from a
                                           %% different sample and does not apply here.

%% the specification ladder, centered velocity throughout (fit.log, placebo_diag.log)
\newcommand{\ladderMinTau}{\ensuremath{+14.62}}   \newcommand{\ladderMinTauSE}{5.47}
\newcommand{\ladderMinPl}{\ensuremath{-7.75}}     \newcommand{\ladderMinPlSE}{1.91}
\newcommand{\ladderMinPlSigmas}{4.0}
\newcommand{\ladderGeoTau}{\ensuremath{-7.62}}    \newcommand{\ladderGeoTauSE}{7.01}
\newcommand{\ladderGeoPl}{\ensuremath{-4.12}}     \newcommand{\ladderGeoPlSE}{1.90}
\newcommand{\ladderFullTau}{\ensuremath{-2.40}}   \newcommand{\ladderFullTauSE}{7.18}
\newcommand{\ladderFullPl}{\ensuremath{-2.66}}    \newcommand{\ladderFullPlSE}{1.71}
\newcommand{\ladderFullPlSigmas}{1.6}
\newcommand{\outTauSigmas}{0.33}           %% E31, tau in standard errors from zero
\newcommand{\outCentroidSigmas}{3.65}      %% E31, standard errors from tau = 23.81 ms
\newcommand{\outInterval}{26.2}            %% ms, evidence centroid minus output time
\newcommand{\outIntervalSE}{7.2}           %% ms, its standard error

%% the sign-flip control: the same fit under three velocity estimators (fit.log arm 1,
%% injection.log arm 3). A forward difference shares the label-center noise with the
%% residual, a backward one carries it with the opposite sign, a centered one omits it.
\newcommand{\sfMinFwd}{\ensuremath{-14.68}}   \newcommand{\sfMinFwdSE}{4.99}
\newcommand{\sfMinBack}{\ensuremath{+45.65}}  \newcommand{\sfMinBackSE}{5.59}
\newcommand{\sfMinCen}{\ensuremath{+17.51}}   \newcommand{\sfMinCenSE}{5.06}
\newcommand{\sfMinPlFwd}{\ensuremath{-6.53}} \newcommand{\sfMinPlFwdSE}{1.68}
\newcommand{\sfMinPlBack}{\ensuremath{-6.15}}\newcommand{\sfMinPlBackSE}{1.71}
\newcommand{\sfMinPlCen}{\ensuremath{-7.76}} \newcommand{\sfMinPlCenSE}{1.92}
\newcommand{\sfFullFwd}{\ensuremath{-20.60}}  \newcommand{\sfFullFwdSE}{6.36}
\newcommand{\sfFullBack}{\ensuremath{+14.38}} \newcommand{\sfFullBackSE}{5.97}
\newcommand{\sfFullPlFwd}{\ensuremath{-1.99}} \newcommand{\sfFullPlFwdSE}{1.39}
\newcommand{\sfFullPlBack}{\ensuremath{-2.10}}\newcommand{\sfFullPlBackSE}{1.50}
\newcommand{\sfSpread}{35.0}               %% ms, forward to backward under full controls
\newcommand{\sfMidpoint}{\ensuremath{-3.11}}  %% ms, their midpoint, against -2.40 centered

%% the placebo diagnosis (placebo_diag.log): splitting by the sign of the horizontal
%% velocity flips the rotated-velocity coefficient, which a rotational term cannot do.
\newcommand{\plRight}{\ensuremath{-27.99}} \newcommand{\plRightSE}{3.82}
\newcommand{\plRightSigmas}{7.3}
\newcommand{\plLeft}{\ensuremath{+21.84}}  \newcommand{\plLeftSE}{5.47}
\newcommand{\plLeftSigmas}{4.0}

%% injection (injection.log arms 1 and 2): the controls do not absorb a lag.
\newcommand{\injRecovery}{1.000}           %% recovered fraction at 10, 25, 50, 100 ms
\newcommand{\injGeoGamma}{0.05}            %% px per px, the geometric displacement injected
\newcommand{\injGeoMin}{\ensuremath{+0.63}}%% ms, how far it moves the minimal specification
\newcommand{\injGeoFull}{0.01}             %% ms, how far it moves the full one, below this

%% the speed-floor arm (placebo_diag.log section 6), full controls throughout
\newcommand{\floorFiveTau}{\ensuremath{-2.92}}    \newcommand{\floorFiveTauSE}{9.05}
\newcommand{\floorTwentyTau}{\ensuremath{-20.24}} \newcommand{\floorTwentyTauSE}{14.78}

%% ===== E28 measured 2026-09-07 (experiments/e28_gen1_label_noise): the label-center
%% noise of Gen1, measured on Gen1 by E13's difference ladder. E13's 0.4974 px is
%% DSEC-Det's and does not transfer.
\newcommand{\labelSigmaGen}{0.6515}        %% px per axis, E28, sigma_total
\newcommand{\labelSigmaGenAnn}{0.6353}     %% px, its annotation part
\newcommand{\labelSigmaGenQuant}{0.1443}   %% px, its quantization part, 0.5 px lattice
\newcommand{\labelTracksGen}{691}          %% E28, tracks of >= 12 labels on uniform 250/500/1000-ms grids

%% ===== E29 measured 2026-09-07 (experiments/e29_synthetic_signflip): the
%% forward-difference artefact, calibrated on synthetic tracks with tau_true = 0.
\newcommand{\synFwd}{\ensuremath{-5.64}}   %% ms, forward estimator, no lag present
\newcommand{\synBack}{\ensuremath{+5.84}}  %% ms, backward estimator, no lag present
\newcommand{\synCen}{\ensuremath{+0.00}}   %% ms, centered estimator, no lag present
\newcommand{\synSum}{\ensuremath{+0.20}}   %% ms, forward plus backward
\newcommand{\synPred}{\ensuremath{-5.63}}  %% ms, -Var(delta)/dt / Var(v), closed form
\newcommand{\synZero}{\ensuremath{+0.02}}  %% ms, all three with the label noise removed
\newcommand{\synRecover}{0.92}             %% fraction of an injected lag recovered
\newcommand{\synLambda}{0.933}             %% the errors-in-variables prediction
\newcommand{\synTracks}{2000}              %% synthetic tracks

%% displacement the interval amounts to, at speeds that occur in the split (E32)
\newcommand{\outSideMedian}{41.7}          %% px, median box side of the matches
\newcommand{\outSpeedNinety}{25.7}         %% px/s, 90th percentile label speed
\newcommand{\outSpeedNinetynine}{54.2}     %% px/s, 99th percentile label speed
\newcommand{\outPxNinety}{0.61}            %% px, displacement over 23.81 ms at p90
\newcommand{\outPxNinetynine}{1.29}        %% px, the same at p99
\newcommand{\outStaticPct}{22}             %% percent of matches below 2 px/s

%% ===== E25 measured 2026-09-04 (experiments/e25_six_sequences): the mains
%% signature on ALL SIX DSEC train sequences whose exposure is pinned at 14996 us.
%% Per labelled box, the modulation depth m = 2|C| and the Rayleigh statistic
%% Z = N|C|^2 at 100 Hz, with an off-frequency arm at 137 Hz. For these finite
%% exposure windows the per-box analytic phase null is not used; E38 supplies
%% an empirical off-frequency reference.
\newcommand{\sixSeqs}{six}                 %% E25, ceiling sequences measured
\newcommand{\sixBoxes}{16\,516}            %% E25, labelled boxes over the six
\newcommand{\sixZmedian}{114.58}           %% E25, median over sequences of the median box Z at 100 Hz
\newcommand{\sixZctrl}{4.01}               %% E25, the same at 137 Hz
\newcommand{\sixLockedLoPct}{95.4}         %% percent, E38b: the LOWEST per-sequence fraction of boxes whose local excess
                                           %% at 100 Hz exceeds the 99th percentile of the
                                           %% same quantity at 137 Hz. This replaces E25's
                                           %% p<1e-3 fraction, withdrawn with the analytic
                                           %% per-box null (experiments/e38_empirical_null).
\newcommand{\sixLockedHiPct}{100.0}        %% percent, E38b, highest (same construction)
\newcommand{\sixModLo}{0.428}              %% E25, lowest per-sequence median modulation depth
\newcommand{\sixModHi}{0.489}              %% E25, highest
%% per-sequence rows of Table 2, in the order printed
\newcommand{\sixBoxesA}{2914}   \newcommand{\sixModA}{0.428}
\newcommand{\sixZA}{117.6}      \newcommand{\sixZcA}{3.15}  \newcommand{\sixLockA}{100.00}
\newcommand{\sixBoxesB}{2134}   \newcommand{\sixModB}{0.440}
\newcommand{\sixZB}{99.1}       \newcommand{\sixZcB}{4.69}  \newcommand{\sixLockB}{99.63}
\newcommand{\sixBoxesC}{619}    \newcommand{\sixModC}{0.445}
\newcommand{\sixZC}{68.4}       \newcommand{\sixZcC}{2.75}  \newcommand{\sixLockC}{99.68}
\newcommand{\sixBoxesD}{241}    \newcommand{\sixModD}{0.489}
\newcommand{\sixZD}{310.9}      \newcommand{\sixZcD}{25.85} \newcommand{\sixLockD}{95.44}
\newcommand{\sixBoxesE}{6997}   \newcommand{\sixModE}{0.466}
\newcommand{\sixZE}{251.5}      \newcommand{\sixZcE}{7.82}  \newcommand{\sixLockE}{99.23}
\newcommand{\sixBoxesF}{3611}   \newcommand{\sixModF}{0.438}
\newcommand{\sixZF}{111.6}      \newcommand{\sixZcF}{3.33}  \newcommand{\sixLockF}{100.00}
\newcommand{\sixMinEvents}{400}            %% E25, events per box required

%% ===== E26 measured 2026-09-04 (experiments/e26_egomotion, paired.json): the
%% control that retracted E20. Each box's fitted gradient is compared against the
%% gradient of an equal-area annulus of its own surrounding background, aligned
%% with the same label velocity, and the difference is taken per box so that a
%% cause common to the box and its surroundings cancels.
\newcommand{\egoRingBetter}{four}          %% E26, sequences where the ring aligns better
\newcommand{\egoCoreNeg}{five}             %% E26, sequences where the eroded core goes negative
\newcommand{\egoCoreWorst}{\ensuremath{-0.178}}  %% E26, paired core minus ring, zurich_city_01_a
\newcommand{\egoCoreWorstSigmas}{6.2}      %% E26, that difference in standard errors

%% --- E02, provenance of the window facts (src/e02_verify_window.py) ---------
\newcommand{\rvtSourceLine}{405}           %% line of preprocess_dataset.py that builds
                                           %% ev_repr_timestamps_us_end backwards from labels

%% --- the two temporal terms on the pixel scale of the split they were measured on

%% ===== E32b measured 2026-09-07 (experiments/e27_rows/bootstrap.log): a cluster
%% bootstrap over sequences for the reported specification, fitted by
%% Frisch-Waugh-Lovell (the group dummies are absorbed by within-cell demeaning,
%% which reproduces the full design's tau to 0.007 ms).
\newcommand{\outBootLo}{\ensuremath{-17.0}}   %% ms, 2.5th percentile
\newcommand{\outBootHi}{\ensuremath{+11.5}}   %% ms, 97.5th percentile
\newcommand{\outBootDraws}{4000}              %% resamples
\newcommand{\outBootBeyond}{3}                %% draws at or beyond 23.81 ms, of 4000

%% ===== E33 measured 2026-09-07 (experiments/e27_rows/saturate.log): the reported
%% specification pairs each geometric control with one image axis. This enters all
%% four on BOTH axes, drops each in turn, and restricts the sample, so that the
%% spread of tau across specifications is reported rather than one of them chosen.
\newcommand{\satTau}{\ensuremath{-3.62}}   \newcommand{\satTauSE}{7.41}
\newcommand{\satPl}{\ensuremath{-1.59}}    \newcommand{\satPlSE}{1.62}
\newcommand{\satLo}{\ensuremath{-20.2}}    %% ms, lowest tau over EVERY specification
                                           %% carrying geometric controls, including the
                                           %% 20 px/s floor of placebo_diag.log, the
                                           %% saturate.log sweep alone gives -9.8.
\newcommand{\satHi}{\ensuremath{+11.1}}    %% ms, highest
\newcommand{\satTight}{\ensuremath{+0.49}} \newcommand{\satTightSE}{4.67}
\newcommand{\satTightIoU}{0.7}             %% the overlap floor giving the tightest fit

%% ===== E34 measured 2026-09-07 (experiments/e34_influence_robust/result.json): is the
%% influence centroid an artefact of occluding a bin with zeros? Three fills and one
%% instrument that does not occlude at all. The recurrent state is held fixed within
%% every sample, so none of these measures state perturbation.
\newcommand{\occZero}{\ensuremath{-23.81}}   %% ms, zeros, what E17 did
\newcommand{\occMean}{\ensuremath{-26.35}}   %% ms, the bin set to its dataset mean
\newcommand{\occSwap}{\ensuremath{-25.01}}   %% ms, the bin taken from an unrelated sample
\newcommand{\occGrad}{\ensuremath{-22.35}}   %% ms, d||out||/d(bin), no occlusion at all
\newcommand{\occSpread}{4.01}                 %% ms, the range across all four instruments
\newcommand{\occBootSE}{0.054}                %% ms, bootstrap SE of the zero-fill centroid
\newcommand{\occGradCV}{0.191}               %% CV over bins of the gradient instrument

%% ===== E35 measured 2026-09-07 (experiments/e27_rows/ratio.log): the nearest prior
%% estimator, the per-object ratio, fitted on exactly the rows the regression uses.
\newcommand{\ratioMean}{\ensuremath{+24.22}} \newcommand{\ratioMeanSE}{27.73}
\newcommand{\ratioIQR}{430}                  %% ms, its interquartile range over all rows
\newcommand{\ratioSD}{4170}                  %% ms, its sample standard deviation
\newcommand{\ratioFloor}{\ensuremath{+13.34}}\newcommand{\ratioFloorSE}{6.17}
\newcommand{\ratioFloorSpeed}{15}            %% px/s, the floor that arm imposes

%% plot and procedure constants, so that no numeral appears in prose
\newcommand{\bootCI}{95}                     %% percent, the bootstrap interval reported
\newcommand{\injLevels}{10, 25, 50 and 100}  %% ms, the injected lags of injection.log
\newcommand{\flickerDayPeakHz}{103.0}        %% Hz, where the day sequence's in-band peak sits
                                             %% (e10_flicker/verify.json, DAY top_line_Hz)

%% errors-in-variables on the reported rows (experiments/e27_rows/fit.log, section 5,
%% re-run with SIGMA_C from E28 rather than the E13 default)
\newcommand{\eivVarObs}{110.6}             %% px^2/s^2, Var of the centered velocity
\newcommand{\eivVarNoise}{1.0}             %% px^2/s^2, row-wise mean 2 sigma^2 / (dt_f+dt_b)^2
\newcommand{\eivLambda}{0.991}             %% the attenuation factor on these rows
\newcommand{\eivShift}{0.2}                %% ms, conservative rounded upper bound on correction shift

%% ===== E38 measured 2026-09-07 (experiments/e38_empirical_null): the per-box Rayleigh
%% null, measured rather than assumed. RETRACTION: the per-box p-values and the
%% "99.0-100.0 % of boxes locked at p<1e-3" figure are withdrawn. Pooled over the six
%% ceiling sequences the OFF-FREQUENCY per-box Z has a median of 38.39 against the
%% analytic 0.693, an inflation of 55, because event times inside one box during a 15 ms
%% exposure arrive in bursts. The per-PIXEL arm is unaffected: there the 137 Hz control
%% returns 0.581 against 0.693, so the analytic null holds where it was verified.
\newcommand{\nullOffMed}{38.4}             %% pooled per-box Z at off frequencies
\newcommand{\nullInflation}{55}            %% times the analytic median
\newcommand{\nullOffFreqs}{eight}             %% off-frequencies swept

%% arm 2 (local_excess.log): R = Z(100 Hz) / the same box's own local background, which
%% needs no null at all. Grid 60-300 Hz at 1 Hz, the background band is |f-100| in [6,30]
%% excluding 20 Hz and 100 Hz harmonics.
\newcommand{\localBoxes}{16\,516}          %% boxes pooled over the six sequences
\newcommand{\localLine}{1.87}              %% median R at the line
\newcommand{\localLinePten}{1.62}          %% its tenth percentile
\newcommand{\localLinePninety}{2.18}       %% its ninetieth
\newcommand{\localSham}{0.11}              %% median R at the off-frequency arm
\newcommand{\localShamPninetynine}{1.00}   %% its 99th percentile
\newcommand{\localRatio}{17.5}             %% ratio of the two medians
\newcommand{\localAbovePct}{99.63}         %% percent of boxes above the sham's 99th pct
\newcommand{\localPeakHz}{96}              %% Hz, where the pooled median spectrum peaks
\newcommand{\localPeakOverHigh}{8.5}       %% peak over the 200-300 Hz median
\newcommand{\localPeakWidth}{35}           %% Hz, its half-maximum width
\newcommand{\localTransformLimit}{67}      %% Hz, 1/14996 us, what one window can resolve
\newcommand{\localGridLo}{60}\newcommand{\localGridHi}{300}

%% ===== E37 measured 2026-09-07 (experiments/e37_map): average precision as a function
%% of the instant the ground truth describes. One GPU pass dumps every detection and every
%% ground-truth box of the Gen1 validation split, the sweep displaces the ground truth by
%% delta*v and recomputes mAP. Ground truth without a centered velocity is marked ignore.
\newcommand{\mapFrames}{20\,296}           %% frames evaluated
\newcommand{\mapDets}{86\,788}             %% detections dumped
\newcommand{\mapGT}{40\,698}               %% ground-truth boxes
\newcommand{\mapGTvelPct}{68.7}            %% percent carrying a centered velocity
\newcommand{\mapArgmax}{\ensuremath{-10}}  %% ms, where mAP is maximised
\newcommand{\mapZero}{0.3346}              %% mAP at the label instant
\newcommand{\mapCentroid}{0.3342}          %% mAP with the gt displaced to the centroid
\newcommand{\mapCost}{0.04}                %% mAP points between those two
\newcommand{\mapCostAbs}{0.0004}           %% absolute mAP change
\newcommand{\mapCostPct}{0.12}             %% and as a percentage
\newcommand{\mapSweepSpan}{0.9}            %% mAP points over the whole +-60 ms sweep
\newcommand{\mapSpeedMedian}{4.0}          %% px/s, median label speed of this set
\newcommand{\mapSideMedian}{41.5}          %% px, median box side
\newcommand{\mapPxMedian}{0.10}            %% px, 23.81 ms at that median speed
%% the speed strata (strata.log): span over +-50 ms, and n
\newcommand{\mapSpanSlow}{0.11}   \newcommand{\mapNslow}{20\,154}
\newcommand{\mapSpanMid}{0.33}    \newcommand{\mapNmid}{5\,356}
\newcommand{\mapSpanHigh}{0.78}   \newcommand{\mapNhigh}{2\,096}
\newcommand{\mapSpanTop}{0.37}    \newcommand{\mapNtop}{337}
\newcommand{\mapTopMap}{0.139}    %% the fastest stratum's mAP, against 0.286 below it
\newcommand{\mapHighMap}{0.286}

%% ===== E55 measured 2026-09-15 (experiments/e37_map/stratum_table.json): the speed
%% strata evaluated EXACTLY at the corrected centroid of -23.81 ms (E37c's per-stratum
%% costs sat at the superseded -24.94 ms and on the 5 ms grid), with each stratum's
%% share of the moving boxes, its median speed, and what 23.81 ms of that speed is in
%% pixels. Spans are E37c's, whose grid did not change. The all-moving row reproduces
%% at_centroid.json to inside 5e-5, asserted in the script.
\newcommand{\stShareSlow}{72.1}\newcommand{\stShareMid}{19.2}
\newcommand{\stShareHigh}{7.5}\newcommand{\stShareTop}{1.2}
\newcommand{\stSpdSlow}{2.1}\newcommand{\stSpdMid}{15.0}
\newcommand{\stSpdHigh}{32.3}\newcommand{\stSpdTop}{57.4}
\newcommand{\stPxSlow}{0.05}\newcommand{\stPxMid}{0.36}
\newcommand{\stPxHigh}{0.77}\newcommand{\stPxTop}{1.37}
\newcommand{\stCostSlow}{0.02}\newcommand{\stCostMid}{0.01}
\newcommand{\stCostHigh}{0.08}\newcommand{\stCostTop}{0.00}
\newcommand{\stSpanAll}{0.37}     %% the all-moving span on the -50..+30 ms grid, the
                                  %% 0.9 of \mapSweepSpan is the +-60 ms sweep of E37b
\newcommand{\stNmoving}{27\,943}  %% ground-truth boxes carrying a centered velocity

%% ===== Real-image figures rendered 2026-09-15 (src/make_fig_real.py). fig6: the frame is
%% chosen by rule (>=3 velocity boxes, >=2 at >=20 px/s, largest median speed) from E37's
%% own iteration order, so the drawn boxes are the ones the sweep scored, stats in
%% experiments/e37_map/fig6_stats.json. fig7: densest published ceiling exposure window vs
%% the same instants of the daytime control, stats in
%% experiments/e00_exposure_survey/fig7_stats.json.
\newcommand{\sweepRealFast}{45.2}   %% px/s, the fastest label on the drawn frame
\newcommand{\sweepRealPx}{1.08}     %% px, that speed over the 23.81 ms interval
\newcommand{\ceilVarNight}{3.0}     %% max/min of the 350 us rate, ceiling sequence
\newcommand{\ceilVarDay}{1.1}       %% the same on the daytime control

%% ===== E40 measured 2026-09-07 (experiments/e40_harmonics/result.json): which frequency
%% carries the signature, with no band chosen in advance. E10's "strongest line between 90
%% and 110 Hz" could only ever have found a line in that band. A 20-450 Hz sweep on each
%% whole recording shows the HARMONIC SERIES of a 100 Hz fundamental, strongest at the
%% second harmonic, with nothing at the odd multiples of 50 Hz and nothing in daylight.
\newcommand{\harmOne}{30}                  %% ceiling median, 100 Hz over its own continuum
\newcommand{\harmTwo}{125}                 %% the same at 200 Hz
\newcommand{\harmThree}{15}                %% at 300 Hz
\newcommand{\harmFifty}{1.8}               %% at 50 Hz: no line
\newcommand{\harmOneFifty}{4.5}            %% at 150 Hz: no line
\newcommand{\harmTwoFifty}{1.6}            %% at 250 Hz: no line
\newcommand{\harmDayOne}{0.5}              %% daytime sequence at 100 Hz
\newcommand{\harmDayTwo}{6.2}              %% daytime sequence at 200 Hz
\newcommand{\harmTwoMax}{2518}             %% the strongest, zurich_city_09_a at 200 Hz
\newcommand{\harmTwoMin}{14.7}             %% the weakest ceiling sequence at 200 Hz
\newcommand{\harmSecondStrongest}{four}    %% sequences where 200 Hz is the strongest line
\newcommand{\harmGridLo}{20}\newcommand{\harmGridHi}{450}
\newcommand{\harmSecondHz}{200}\newcommand{\harmThirdHz}{300}
\newcommand{\harmBoxTwo}{0.76}             %% E38b at 200 Hz: the per-box local excess there
                                           %% is WEAKER than at 100 Hz (1.87), because one
                                           %% 14996 us window resolves only 67 Hz

%% ===== E41 measured 2026-09-07 (experiments/e41_permutation_null): the within-frame
%% dispersion against a null that assumes nothing about arrival statistics. Events of all
%% qualifying boxes of a frame are pooled and reassigned at random to boxes, preserving
%% counts, so every temporal structure common to the frame survives and only box identity
%% is destroyed. The analytic null's independence assumption is thereby removed.
\newcommand{\permNull}{78.6}               %% us, the permutation null, median over frames
\newcommand{\permExcess}{191.1}            %% us, the excess over it
\newcommand{\permDraws}{200}               %% draws per frame
\newcommand{\permFrames}{1134}             %% qualifying frames
\newcommand{\permAbovePct}{94.1}           %% percent of frames exceeding their own null
\newcommand{\permRatio}{0.85}              %% permutation over analytic: it is SMALLER
\newcommand{\permRandObs}{339.4}           %% us, random-position boxes, observed
\newcommand{\permRandNull}{113.3}          %% us, their permutation null
\newcommand{\permRandExcess}{320.0}        %% us, their excess

%% retained for provenance only; the finite-window per-box table reports depth
%% descriptively and does not use a random-phase floor.
\newcommand{\modFloorSix}{0.089}           %% sqrt(pi/400)
\newcommand{\fwlAgree}{0.007}              %% ms, FWL against the full design (bootstrap.log)
\newcommand{\mapArgmaxTau}{\ensuremath{+10}}  %% ms, \mapArgmax read on the p-g=-tau v axis
\newcommand{\localRatioSweep}{4.0}         %% the pooled 100 Hz to off-frequency ratio of
                                           %% medians over all eight off frequencies (E38
                                           %% run.log), against 28.6 for the 137 Hz arm alone
\newcommand{\mapVsFit}{1.7}                %% the mAP argmax against the regression, in SEs
                                           %% of the regression: (10-(-2.40))/7.18 = 1.73

%% ===== E42 measured 2026-09-07 (experiments/e42_recurrent_support): the temporal support
%% of a RECURRENT detector, including its history. E17 and E34 hold the recurrent state
%% fixed and perturb the current window only, so their centroid is the CURRENT WINDOW's
%% influence centroid conditional on that state, not the support of the emitted state as
%% Sec. 3.1 defines it. This occludes whole past windows and lets the state propagate.
\newcommand{\histShareTen}{30.0}           %% percent of influence in the current window,
                                           %% 500 ms horizon
\newcommand{\histShareTwenty}{22.7}        %% the same over a 1000 ms horizon
\newcommand{\histCentroidTen}{\ensuremath{-161}}    %% ms, over the FIRST TEN windows of the k=19
                                           %% run, so both horizons come from one sample,
                                           %% the k=9 run began at a different frame and its
                                           %% -166 is not mixed in
\newcommand{\histCentroidTwenty}{\ensuremath{-299}} %% ms, the same over 1000 ms
\newcommand{\histHalfLag}{\ensuremath{-175}}        %% ms, where half the influence is reached
\newcommand{\histResetChange}{0.0867}      %% relative change when the state is zeroed
\newcommand{\histCurrentChange}{0.0231}    %% the same for occluding the whole current window
\newcommand{\histResetRatio}{3.8}          %% the ratio of those two
\newcommand{\histSamples}{576}             %% samples
\newcommand{\histHorizonA}{500}\newcommand{\histHorizonB}{1000}  %% ms
%% the fill control (fill_control.log): the plateau is not an artefact of zero filling
\newcommand{\histDecayZero}{10.3}          %% decay from lag 0 to lag 19, zero fill
\newcommand{\histDecaySwap}{2.2}           %% the same under a swap fill
\newcommand{\histSwapRatioLo}{2.43}       %% swap/zero sensitivity ratio at lag 0
\newcommand{\histSwapRatioHi}{11.47}      %% the same at lag 19
\newcommand{\histTailSwapPct}{46}          %% percent of lag-0 influence still at lag 19

%% stratum edges and band frequencies, so that no numeral enters prose (E37 strata.log,
%% E40 result.json). The p<1e-3 level is a chosen significance and is written as a power.
\newcommand{\mapCutLo}{10}\newcommand{\mapCutMid}{25}\newcommand{\mapCutHi}{50}
\newcommand{\harmOddOneHz}{150}\newcommand{\harmOddTwoHz}{250}
\newcommand{\ladderTrimPct}{5}             %% percent trimmed per tail in E28's ladder

%% ===== E37d measured 2026-09-07 (experiments/e37_map/argmax_bootstrap.log): a frame
%% bootstrap of the mAP sweep's argmax. The argmax does not identify a value, and in no
%% resample does it reach the newest window's centroid.
\newcommand{\mapArgLo}{\ensuremath{-20}}   %% ms, 2.5th percentile of the argmax
\newcommand{\mapArgHi}{\ensuremath{0}}     %% ms, 97.5th percentile
\newcommand{\mapArgBoot}{1000}             %% frame resamples
\newcommand{\mapArgNever}{none}            %% resamples whose argmax reaches -23.81 ms

%% ===== E44 measured 2026-09-07 (experiments/e44_capacity/rvt-{t,s,b}.json): the same
%% predictor-side measurement on the three released Gen1 capacities, which differ only in
%% embed_dim, attention dim_head and FPN depth. Same sequences, representation, warm-up and
%% code for all three, all three load with 0 missing and 0 unexpected keys (E43). The
%% newest-window centroids here use the state entering the step, as E45 does, which is why
%% rvt-t agrees with E45 and not with the superseded E17.
\newcommand{\capSamples}{576}            %% samples per capacity
\newcommand{\capParamsT}{4.41}        \newcommand{\capParamsS}{9.87}
\newcommand{\capParamsB}{18.54}
\newcommand{\capBinT}{\ensuremath{-23.84}}   \newcommand{\capBinS}{\ensuremath{-23.89}}
\newcommand{\capBinB}{\ensuremath{-24.11}}
\newcommand{\capBinSpread}{0.27}   %% ms, range of the three
\newcommand{\capShareT}{22.7}   \newcommand{\capShareS}{20.6}
\newcommand{\capShareB}{19.0}
\newcommand{\capSuppT}{\ensuremath{-299}}  \newcommand{\capSuppS}{\ensuremath{-337}}
\newcommand{\capSuppB}{\ensuremath{-339}}
\newcommand{\capResetT}{3.8}   \newcommand{\capResetB}{3.0}
\newcommand{\capHorizon}{1000}          %% ms, the horizon these are quoted over

%% ===== E47/E48 measured 2026-09-08 (experiments/e47_ssm, experiments/e48_matched_frames):
%% the newest-window bin arm on TWO architectures over the same 12 Gen1 validation
%% sequences and the same frames. SSM-ViT (Zubic et al., CVPR 2024) is a fork of the RVT
%% repository in which each stage's ConvLSTM is replaced by an S5 state space layer, it
%% distributes RVT's own preprocessed Gen1 and builds its window boundaries with the same
%% line of the same script. Both released checkpoints load with 0 missing and 0 unexpected
%% keys. E48 re-runs RVT's three capacities on E47's frames so the two families are
%% measured on identical samples, its rvt-t agrees with E45. SSM-ViT is driven the way its
%% released evaluation drives it, a chunk of windows at a time with the detection head
%% applied per position, and measured at positions carrying half a chunk of history or more
%% (E47b), the one-window-at-a-time convention is withdrawn, see withdrawn_L1/WHY.md.
\newcommand{\archCkpts}{5}            %% released checkpoints measured
\newcommand{\archOther}{4}            %% the ones Fig. 2a draws besides rvt-t
\newcommand{\archSpan}{0.96}       %% ms, range of all their centroids
\newcommand{\archLo}{\ensuremath{-24.72}}   \newcommand{\archHi}{\ensuremath{-23.76}}
\newcommand{\ssmParamsS}{9.68}  \newcommand{\ssmParamsB}{18.19}
\newcommand{\ssmBinS}{\ensuremath{-23.76}}
\newcommand{\ssmBinB}{\ensuremath{-24.72}}
\newcommand{\mfBinT}{\ensuremath{-23.81}}
\newcommand{\mfBinS}{\ensuremath{-23.98}}
\newcommand{\mfBinB}{\ensuremath{-24.13}}
\newcommand{\archSamples}{480}          %% samples per RVT checkpoint
\newcommand{\archSamplesSsm}{288}       %% samples per SSM-ViT checkpoint:
                                           %% fewer, because only chunk positions carrying
                                           %% half a chunk of history or more are used
\newcommand{\archUnifMax}{1.24}       %% ms, the largest distance of any
                                           %% of them from the uniform-weight -25.00
\newcommand{\archUnifPct}{25}        %% that, as a percent of one bin
%% ===== E47c measured 2026-09-08 (experiments/e47_ssm/statepos-*.json): what the
%% recurrent state carried between evaluation chunks does to the output, by position
%% in the chunk, as a relative L2 change in the emitted detection tensor. The control
%% arm occludes the window one position earlier inside the same chunk, which a model
%% that ignored its own history entirely would also return zero for.
\newcommand{\ssmStateEffect}{0.00000}   %% largest over all positions, s5vit-small
\newcommand{\ssmInchunkOne}{0.0276}      %% previous window, position 1
\newcommand{\ssmInchunkLast}{0.0106}     %% the same at the last position
\newcommand{\ssmChunkMeas}{8}          %% windows per chunk in this measurement
\newcommand{\rvtStateEffect}{0.0304}     %% the same arm on rvt-t, last position
\newcommand{\rvtStateEffectOne}{0.0616}  %% and at position 1
\newcommand{\ssmEvalChunk}{21}          %% windows per chunk in the released streaming
                                           %% evaluation: config/dataset/gen1.yaml
                                           %% sequence_length, and sequence_for_streaming.py
                                           %% start_indices step by exactly that

%% ===== E51-E53 measured 2026-09-08 (experiments/e51_ranking): does the interval change a
%% benchmark COMPARISON? Detections for all five released checkpoints were dumped under one
%% protocol (same frames, same confidence, same NMS, ground truth identical across models and
%% identical to E37's), then each model was scored at delta = 0 and again at its own best
%% displacement, in four speed strata. The ordering never changes, and it cannot: in every
%% stratum the largest gain any model draws from alignment is a fraction of the smallest gap
%% between adjacent models, so no refinement of the grid can reorder them.
\newcommand{\rankCkpts}{5}            %% released checkpoints scored
\newcommand{\rankStrata}{four}          %% speed strata examined
\newcommand{\rankFlips}{0}              %% pairwise sign flips, of 10 per stratum
\newcommand{\rankMaxGain}{0.14}      %% mAP points, the largest alignment gain seen
\newcommand{\rankMinGap}{0.11}       %% mAP points, the smallest adjacent gap seen
\newcommand{\rankRatio}{52}          %% percent: worst-case gain over gap within a stratum
\newcommand{\rankArgLo}{\ensuremath{-15}}   %% ms, earliest per-model argmax
\newcommand{\rankArgHi}{\ensuremath{0}}    %% ms, latest per-model argmax
\newcommand{\rankMover}{\texttt{s5vit-base}}  %% the checkpoint whose rank moves most
\newcommand{\rankMoverBest}{1}          %% its best rank across strata
\newcommand{\rankMoverWorst}{5}         %% and its worst

%% ---------------------------------------------------------------------------
%% Resolving power of the metric (E56). What size of timing difference the mAP
%% can separate from its own cluster-resampling noise, over 406 sequences.
%% Emitted by src/e61_macros.py from experiments/e56_resolving/resolving.json.
\newcommand{\seDmapLo}{0.035}
\newcommand{\seDmapHi}{0.058}
\newcommand{\seMapLo}{2.19}
\newcommand{\seMapHi}{2.82}
\newcommand{\resPairedLo}{13.7}
\newcommand{\resPairedHi}{19.8}
\newcommand{\effTauLo}{0.041}
\newcommand{\effTauHi}{0.081}
\newcommand{\twoSeDmapLo}{0.070}
\newcommand{\twoSeDmapHi}{0.116}
\newcommand{\effTauResolved}{1}
\newcommand{\resAbsLo}{153}
\newcommand{\resAbsHi}{174}
\newcommand{\dstarAll}{43.7}
\newcommand{\dstarAllTau}{1.83}
\newcommand{\dstarMid}{11.2}
\newcommand{\dstarMidTau}{0.47}

%% Stability of the published ordering under the benchmark's own noise (E56d),
%% from experiments/e56_resolving/rankse.json.
\newcommand{\rankHoldBox}{42.5}
\newcommand{\rankHoldMov}{80.5}
\newcommand{\rankFlipMax}{0.48}
\newcommand{\rankGapMin}{0.050}
\newcommand{\rankGapMinSE}{0.53}

%% Chunk position in the released SSM-ViT streaming evaluation (E58), from
%% experiments/e58_chunkpos/{allbox,invariant}.json.
\newcommand{\chunkDidBase}{5.32}
\newcommand{\chunkDidBaseSE}{0.53}
\newcommand{\chunkDidBaseZ}{10.0}
\newcommand{\chunkDidSmall}{4.37}
\newcommand{\chunkDidSmallSE}{0.61}
\newcommand{\chunkDidSmallZ}{7.2}
%% Raw full-minus-short mAP difference per checkpoint (E58, all labeled boxes), the
%% quantity main Fig.~2 plots, kept separate from the control-adjusted difference below.
\newcommand{\chunkRawT}{0.01}
\newcommand{\chunkRawS}{0.80}
\newcommand{\chunkRawB}{0.56}
%% The same raw difference net of the placebo trend: for an RVT row the mean of the other
%% two RVT rows, for an SSM row the mean of all three. Positive means history helped.
\newcommand{\chunkNetT}{-0.67}
\newcommand{\chunkNetS}{0.51}
\newcommand{\chunkNetB}{0.16}
\newcommand{\chunkPlaceboMax}{0.67}
\newcommand{\chunkPlaceboZ}{1.56}
\newcommand{\chunkGapBase}{1.75}
\newcommand{\chunkGapSmall}{1.33}
\newcommand{\chunkShortBase}{34.445}
\newcommand{\chunkFullBase}{40.226}
\newcommand{\chunkRelBase}{38.9}
\newcommand{\chunkConfOne}{0.463}
\newcommand{\chunkConfFull}{0.574}
\newcommand{\chunkConfPct}{23.8}

%% Which summary of the label composition sets the resolution, and what buying it
%% costs (E56e), from experiments/e56_resolving/collapse.json and composition.json.
\newcommand{\collapseCVraw}{0.20}
\newcommand{\collapseCVmed}{0.60}
\newcommand{\collapseQbest}{85}
\newcommand{\collapseCVbest}{0.29}
\newcommand{\collapseCV}{0.042}
\newcommand{\collapseD}{0.0628}
\newcommand{\collapseDsd}{0.0029}
\newcommand{\collapseIoU}{0.882}
\newcommand{\collapseCVten}{0.10}
\newcommand{\collapseCVqtr}{0.28}
\newcommand{\tailGainTop}{2.46}
\newcommand{\tailSeTop}{9.11}
\newcommand{\tailNetTop}{0.27}
\newcommand{\tailNetMid}{0.42}
\newcommand{\stNtop}{337}
\newcommand{\rGenAll}{0.095}
\newcommand{\rDsecTrain}{1.31}
\newcommand{\ddRatioMed}{13.9}
\newcommand{\ddBoxesTrain}{290\,163}
\newcommand{\qGenAll}{0.70}
\newcommand{\qDsecTrain}{5.59}
\newcommand{\qDsecRatio}{7.97}
\newcommand{\qDsecNet}{25.7}

%% The recurrent influence tail (E59), from experiments/e42_recurrent_support/tail.json.
\newcommand{\tailExp}{0.645}
\newcommand{\tailExpSE}{0.005}
\newcommand{\tailExpLo}{0.636}
\newcommand{\tailExpHi}{0.655}
\newcommand{\tailRatio}{1.864}
\newcommand{\tailRatioLo}{1.851}
\newcommand{\tailRatioHi}{1.876}
\newcommand{\tailOutside}{77.3}
\newcommand{\tailRmsPower}{2.5}
\newcommand{\tailRmsExp}{23.8}
\newcommand{\tailRmsRatio}{9.4}

%% Per-model cells of the supplement's chunk-position table (E58, all labelled boxes).
%% Three decimals: the smallest difference the table reports is 0.014 pt, so at one
%% decimal the short and full columns would not reproduce the raw column they are
%% differenced into. At three they close against the printed 2-decimal raw values.
\newcommand{\chunkShortSmall}{31.578}
\newcommand{\chunkFullSmall}{36.410}
\newcommand{\chunkRelSmall}{35.6}
\newcommand{\chunkShortT}{38.180}
\newcommand{\chunkFullT}{38.194}
\newcommand{\chunkShortS}{38.371}
\newcommand{\chunkFullS}{39.167}
\newcommand{\chunkShortB}{41.561}
\newcommand{\chunkFullB}{42.126}
\newcommand{\chunkDidT}{0.67}
\newcommand{\chunkSeT}{0.43}
\newcommand{\chunkDidS}{-0.51}
\newcommand{\chunkSeS}{0.48}
\newcommand{\chunkDidB}{-0.16}
\newcommand{\chunkSeB}{0.55}
\newcommand{\aFitShort}{0.628}
% E58e chunk-position curve (src/e58e_curve.py, src/e58f_fig.py)
\newcommand{\chunkRiseBase}{5.78}
\newcommand{\chunkRiseSmall}{4.83}
\newcommand{\chunkRiseRvt}{0.80}
% E58g position-permutation null (src/e58g_perm.py)
\newcommand{\chunkPermSd}{0.55}
\newcommand{\chunkPermZBase}{9.8}
\newcommand{\chunkPermZSmall}{7.6}
\newcommand{\chunkPermPlaceboZ}{1.52}
\newcommand{\chunkPermP}{0.005}
\newcommand{\chunkPermB}{200}
% E62 association-rule sweep (src/e62_association.py)
\newcommand{\assocRules}{14}
\newcommand{\assocLo}{-8.96}
\newcommand{\assocHi}{-0.90}
\newcommand{\assocSpread}{8.06}
\newcommand{\assocSE}{7.12}
\newcommand{\assocGate}{-2.96}
\newcommand{\assocGateSE}{7.12}
\newcommand{\assocFree}{-5.58}
\newcommand{\assocFreeSE}{7.86}
\newcommand{\assocFreeN}{22\,895}
\newcommand{\assocGateN}{22\,518}
%% the rebuild's own gate against the reported table: |(-2.963)-(-2.40)|/7.123
\newcommand{\assocGateVsTable}{0.08}
\newcommand{\assocDelta}{0.37}
\newcommand{\assocMinLo}{12.01}
\newcommand{\assocMinHi}{16.35}
% E63 checkpoint-port parity (src/e63_port_parity.py)
\newcommand{\portEvalGap}{0.101}
\newcommand{\portParamDev}{0.051}
\newcommand{\portSeqs}{429}
\newcommand{\portApT}{38.95}
\newcommand{\portApOwnT}{38.90}
\newcommand{\portApFifty}{69.7}
% E64 sequence-clustered bootstrap of the newest-window profile (src/e64_seqboot.py)
\newcommand{\occSeqs}{12}
\newcommand{\occSeqCiLo}{-24.07}
\newcommand{\occSeqCiHi}{-23.56}
\newcommand{\occSeqSE}{0.13}
\newcommand{\occSeqWidest}{0.51}
\newcommand{\occSeqLoo}{0.08}
\newcommand{\occPerSeqLo}{-24.63}
\newcommand{\occPerSeqHi}{-23.20}

%% E60, the same-frame chunk-boundary intervention: every chunk boundary moved
%% \ssmPairedShift{} windows later, so the frames the release gave one to four windows of
%% history are re-scored with seventeen to twenty. Within model, within frame, identical
%% ground truth and weights; no control model and no parallel-trends assumption.
%% Cluster bootstrap over the same 406 sequences E56 and E58 use.
%% Both released SSM checkpoints are shifted. Suffixes follow \chunkDid*: Small is
%% s5vit-small, Base is s5vit-base, the checkpoint the differenced headline is quoted on.
\newcommand{\ssmPairedN}{3\,574}
\newcommand{\ssmPairedShift}{5}
\newcommand{\ssmPairedMapVelBase}{5.23}
\newcommand{\ssmPairedMapVelBaseSE}{0.49}
\newcommand{\ssmPairedMapVelSmall}{4.81}
\newcommand{\ssmPairedMapVelSmallSE}{0.52}
\newcommand{\ssmPairedMapAllBase}{4.20}
\newcommand{\ssmPairedMapAllSmall}{4.10}
\newcommand{\ssmPairedConfBase}{0.0579}
\newcommand{\ssmPairedConfBaseSE}{0.0048}
\newcommand{\ssmPairedConfSmall}{0.0509}
\newcommand{\ssmPairedConfSmallSE}{0.0035}
\newcommand{\ssmPairedDpfBase}{3.49}
\newcommand{\ssmPairedDpfSmall}{5.09}

%% Remaining E60 rows, for the supplement's paired table. Both channels move: the boxes
%% that match a label are scored higher, and far fewer boxes are emitted in total.
\newcommand{\ssmPairedB}{300}
\newcommand{\ssmPairedMapAllBaseSE}{0.36}
\newcommand{\ssmPairedMapAllSmallSE}{0.32}
\newcommand{\ssmPairedConfTpBase}{0.0409}
\newcommand{\ssmPairedConfTpBaseSE}{0.0040}
\newcommand{\ssmPairedConfTpSmall}{0.0519}
\newcommand{\ssmPairedConfTpSmallSE}{0.0043}
\newcommand{\ssmPairedDpfBaseSE}{0.25}
\newcommand{\ssmPairedDpfSmallSE}{0.26}
\newcommand{\ssmPairedMapVelBaseZ}{10.7}
\newcommand{\ssmPairedMapVelSmallZ}{9.3}
\newcommand{\ssmPairedConfBaseZ}{12.0}
\newcommand{\ssmPairedConfSmallZ}{14.6}
%% The qualitative pair in Fig. 2, and the population it was chosen from. The four
%% count macros record what the panels themselves draw, so the drawn text and this
%% file are checked against one source. The frame is the
%% clearest single instance, so the per-frame averages are quoted beside it in the caption.
\newcommand{\pairedFigLabels}{4}
\newcommand{\pairedFigMatchShort}{0}
\newcommand{\pairedFigMatchFull}{3}
\newcommand{\pairedFigUnmatchShort}{2}
\newcommand{\pairedFigUnmatchFull}{0}
\newcommand{\pairedFigShow}{0.3}
\newcommand{\pairedFigMeanMatchShort}{1.40}
\newcommand{\pairedFigMeanMatchFull}{1.53}
\newcommand{\pairedFigMoreMatched}{13.1}
\newcommand{\pairedFigFewerMatched}{2.7}


\newcommand{\micro}{\ensuremath{\mu}}
\newcommand{\figasset}[2]{%
  \IfFileExists{#1}{\includegraphics[width=#2]{#1}}%
  {\fbox{\begin{minipage}[c][3.4cm][c]{\dimexpr#2-2\fboxsep-2\fboxrule\relax}%
   \centering\small\texttt{#1}\\[2pt]\textit{asset not yet rendered}%
   \end{minipage}}}}
\newcommand{\eff}{\ensuremath{\hat{\tau}}}

\title{One Timestamp Does Not Identify Temporal Support\\
in Event-Detection Benchmarks\\[0.3em]
\large Supplementary Material}
\author{Anonymous CVPR submission}

\begin{document}
\maketitle

\noindent
This document holds the material referred to from the main paper: the
errors-in-variables treatment of the velocity regressor, the protocol behind the
occlusion measurement, the data paths of the DSEC measurements, and additional
implementation and control analyses. The RVT source revision is
\texttt{b80f5683a6e2d5de65d4bde8105d796ccb50dbb1}, and the SSM-ViT source revision is
\texttt{7c871b55a0c5f00673c2c3975f6b6a1c5adbab88}. Section numbers below refer to this document. References labeled
main Sec.~$n$ refer to the main paper.

\section{The velocity regressor}
\label{sup:denominator}

The regression of main Eq.~2 uses the finite-differenced label velocity
$\tilde v=v+\eta$, where $v$ is latent true velocity and $\eta$ is label noise.
Three error mechanisms enter, and they behave differently. The shared-noise term
determines the estimator choice.

\paragraph{Classical dilution.} Let the label center carry noise of standard
deviation $\sigma_{c}$ per frame and let the speed be formed by a central
difference over $\pm\Delta$. The differenced speed then carries noise with variance

\begin{equation}
\sigma_{\eta}^{2}=\frac{2\sigma_{c}^{2}}{(\Delta_f+\Delta_b)^{2}},
\label{eq:dilution}
\end{equation}

and $\tau$ is attenuated toward zero by
$\lambda_{\mathrm{att}}=\mathrm{Var}(v)/(\mathrm{Var}(v)+\sigma_{\eta}^{2})$.
The noise is the label noise of the dataset the fit runs on, so it is measured
there: Sec.~\ref{sup:labelnoise} solves for $\sigma_{c}=\labelSigmaGen$\,px on
Gen1 itself, and Eq.~\ref{eq:dilution} is evaluated row by row rather than
assuming a fixed cadence. On the
reported rows $\mathrm{Var}(\eta)$ is \eivVarNoise{} against a
$\mathrm{Var}(\tilde v)$ of \eivVarObs{} in px$^{2}$/s$^{2}$, so
$\lambda_{\mathrm{att}}=\eivLambda{}$ and the correction moves $\tau$ by less
than \eivShift\,ms. This marginal errors-in-variables calculation indicates
small attenuation. The synthetic arm of Sec.~\ref{sup:signflip} injects a known
lag under a much larger noise and recovers \synRecover{} of it against a
predicted \synLambda{}, so the correction is calibrated as well as small.
The estimated attenuation is only 0.9\%, moving $\tau$ by less than 0.2\,ms,
and therefore cannot explain the near-zero coefficient reported here.

\paragraph{Smoothing bias.} On the uniform grid used for the label-noise
calibration, a central difference returns the mean velocity over $2\Delta$
rather than the instantaneous velocity. Expanding the trajectory about the
query instant,

\begin{equation}
\frac{x(t+\Delta)-x(t-\Delta)}{2\Delta}
=\dot{x}(t)+\frac{\Delta^{2}}{6}\dddot{x}(t)+O(\Delta^{4}),
\end{equation}

so the regressor carries a systematic term of order $\Delta^{2}\dddot{x}/6$. That
is a bias in the regressor, not noise in it, and no errors-in-variables
correction removes it. It is made small by making $\Delta$ small, which
Eq.~\ref{eq:dilution} says makes the dilution worse.
The reported regression rows use the local secant
$[x(t+\Delta_f)-x(t-\Delta_b)]/(\Delta_f+\Delta_b)$. Their median
$|\Delta_f-\Delta_b|$ is zero, so the first-order acceleration term is absent
in the typical row. The general expansion includes
$(\Delta_f-\Delta_b)\ddot{x}/2$ and is covered by the stated local-secant
interpretation.

\paragraph{The shared-noise term.} The first two mechanisms are properties of the
regressor alone. The third couples the regressor to the response. A forward
difference $(c^{*}_{k+1}-c^{*}_{k})/\Delta$ carries the label noise $\delta_{k}$ of
the current center with weight $-1/\Delta$, and the residual
$\hat{c}-c^{*}_{k}$ carries the same $\delta_{k}$ with weight $-1$, so

\begin{equation}
\operatorname{Cov}\!\left(\hat{c}-c^{*}_{k},\,v_{f}\right)=\frac{\operatorname{Var}(\delta)}{\Delta}
\label{eq:shared}
\end{equation}

\noindent even at $\tau=0$, and the induced slope is
$\operatorname{Var}(\delta)/\Delta\operatorname{Var}(v)$. A backward difference
carries $\delta_{k}$ with the opposite sign. The centered difference omits
$c^{*}_{k}$ and removes the direct same-frame shared-noise term. Residual bias
from temporally correlated annotation errors is not identified by this
construction, which is why the forward and backward fits are reported beside it
as a control.

\paragraph{Instrumental-Variable Alternative.} An instrumented
estimate would replace the current speed by a velocity finite differenced from a
strictly disjoint earlier label pair, on the argument that disjointness implies
independent errors. DSEC provides a counterexample to assuming that disjoint
label errors are generally independent. The third-difference kernel $(1,-3,3,-1)$ forces
$\rho(1)=\labelAcfTheory$ for white noise, the same code returns
$\labelAcfWhite$ on white noise and $\labelAcfMeasured$ on the DSEC-Det labels,
and pooled-ratio bias, heteroscedastic segments and heavy tails each returned the
white-noise value through the identical pooling code, so the departure is a
property of the labels. That measurement does not identify what the correlation
structure is, beyond excluding a first-order autoregressive form. The paper does
not report an instrumental estimate.

\paragraph{The construction term.} A benchmark whose high-rate labels are
constructed between anchors adds a term to $e_{\parallel}$, where
$e_{\parallel}=\langle\hat c-c^{*},\hat u\rangle$, that has nothing to do
with the predictor. The chord of a curved image trajectory runs ahead of the
truth throughout the interval and meets it at the anchors, which produces a
signed, frame-locked, resetting residual for a predictor with no lag. Main
Sec.~4.1 shows that exact adjacent-frame linear interpolation is inconsistent
with the observed DSEC-Det label trajectories as the sole construction mechanism,
and the Gen1 measurements of main Sec.~3 are taken at label times.

\section{The influence profile}
\label{sup:influence}

Main Sec.~3.3 measures the occlusion-sensitivity profile of the released
\texttt{rvt-t} checkpoint. The stored input is a stacked histogram of
\rvtBins{} bins of \rvtBinWidth\,ms covering the \rvtWindow\,ms preceding the
query time. For each sample the recurrent state is advanced over
\rvtWarmSteps{} preceding steps so that the network is not read out of a zero
state, the reference output is recorded, and then each bin in turn is zeroed and
the output recorded again. A bin's influence is the magnitude of the change its
occlusion produces in the emitted detection tensor, relative to its norm,
averaged over
\rvtInfluenceSamples{} samples drawn from the first twelve lexicographically
ordered released Gen1 validation sequences, using the first forty post-warm-up
samples from each sequence. This deterministic rule gives 480 samples. The
measured occlusion-sensitivity centroid is the weighted average
$\sum_{k}s_{k}c_{k}/\sum_{k}s_{k}$ over the ten bin centers $c_{k}$, where
$s_k$ denotes the measured non-negative bin-ablation sensitivity. The occluded pass starts from the recurrent
state \emph{entering} the step, not from the state the reference pass returned. The two
passes therefore share one history and differ only in the occluded bin. If the
returned state is used instead, the passes are one step out of phase and the centroid
moves by \rvtCentroidDiff\,ms.

For the recurrent-history measurement, the same deterministic sampling rule uses
\histSamples{} samples, the first forty-eight post-warm-up samples from each of the
same twelve validation sequences. The current window and the preceding nine windows
form the \histHorizonA\,ms horizon, and the current window and preceding nineteen windows
form the \histHorizonB\,ms horizon. For each lag, the corresponding whole input window
is zeroed while the saved reference history before that lag is restored, and the state
is propagated through the intervening windows to the target step. The relative $L_{2}$
change in the target detection tensor is averaged over samples to give a non-negative
window sensitivity $q_{\ell}$. The horizon centroid is
\[
c_{q,H}=\frac{\sum_{\ell=0}^{L_H-1}q_{\ell}[-(25+50\ell)\,\mathrm{ms}]}
{\sum_{\ell=0}^{L_H-1}q_{\ell}},
\]
where $L_H$ is the number of windows in horizon $H$, and the newest-window share is
$q_{0}/\sum_{\ell}q_{\ell}$. The same saved reference
history is used for the reference and ablated passes at each lag.

Two properties of that measurement bound it. Occlusion reports the sensitivity of
the output to a bin. The ablation profile is an operational input-sensitivity
measure and is not, in general, a decomposition of temporal weights. For a
nonlinear predictor it is a local rather than a global description. And a single-bin occlusion
leaves the remaining bins in place, so the measurement is a partial derivative
rather than a decomposition of the output.

Masking the leading bins of the stored input is not a route to the same quantity
at a shorter support. For a fixed network of influence weights $a_{k}$ on bin
centers $c_{k}$, let a constant-velocity trajectory be
$p(t)=p_{0}+vt$. The conditional expectation after masking the leading $j$ bins
is $A_{j}p_{0}+C_{j}v$ with $A_{j}=\sum_{S}a_{k}$ and $C_{j}=\sum_{S}a_{k}c_{k}$
over the retained set $S$. Reading a time off that requires the normalization
$C_{j}/A_{j}$, which the network need not perform, and the result equals the
retained-mass centroid only for non-negative weights.

\section{The label-center noise of Gen1}
\label{sup:labelnoise}

The errors-in-variables treatment of Sec.~1 and the calibration of the velocity
estimator both need the dispersion of the Gen1 label center about a smooth
trajectory. The value solved for elsewhere in this work, \sigmaCtotal\,px, is
DSEC-Det's, and does not transfer to a different annotation effort on a different
sensor, so it is solved again on Gen1.

For a track center series $x_i=f(t_i)+e_i$ on a uniform grid, the $k$-th
difference has $\operatorname{Var}(D_k)=\binom{2k}{k}\sigma^{2}$ plus a motion
term of order $k$, so $\operatorname{Var}(D_k)/\binom{2k}{k}$ falls toward
$\sigma^{2}$ from above. Halving the sampling rate multiplies the motion term by
$4^{k}$ and leaves $\sigma^{2}$ untouched, giving a second equation in the same
two unknowns, $\sigma^{2}=(4^{k}V_{\text{full}}-V_{\text{half}})/(4^{k}-1)$.
Tracks are rebuilt by the same-class overlap linking the output-time
measurement uses, giving \labelTracksGen{} tracks of at least twelve labels on
uniform annotation grids with cadences of 250, 500, or 1000\,ms, corresponding
to 4, 2, or 1\,Hz. The reported regression rows instead use the two adjacent
label neighbors with variable intervals. Their row-wise noise variance is
$2\sigma_c^2/(\Delta_f+\Delta_b)^2$, rather than a fixed-cadence approximation.
Variances are trimmed at \ladderTrimPct\,\%, because all four Gen1 box
coordinates are stored as integers and a median absolute deviation on
lattice-valued data locks onto discrete values.

The two-rate elimination agrees with the full-rate ladder to four decimals from
$k=5$ onward, which is the check that the motion term is negligible rather than
assumed away. The ladder extrapolates to $\sigma=\labelSigmaGen$\,px, of which
\labelSigmaGenQuant\,px is the quantization of a \sigmaClattice\,px center lattice, leaving
\labelSigmaGenAnn\,px of annotation dispersion.

\section{The sign-flip calibration}
\label{sup:signflip}

Equation~\ref{eq:shared} predicts an artifact whose size is set by the label noise
and the observed velocity variance, so it is calibrated before the estimator is
used. \synTracks{} synthetic tracks of fourteen labels on the Gen1 grid are
generated with quadratic motion. The synthetic experiment calibrates the
shared-noise artifact and attenuation under quadratic motion. The simulated detector reports the true position
at $t-\tau_{\text{true}}$ with its own independent error, and the label centers
are corrupted at \labelSigmaGen\,px.

With $\tau_{\text{true}}=0$ the forward estimator returns \synFwd\,ms and the
backward one \synBack\,ms, summing to \synSum\,ms, against the closed form
\synPred\,ms, the centered estimator returns \synCen\,ms. Setting the label
noise to zero returns all three to \synZero\,ms, and the placebo stays at its
null throughout, so it does not absorb the artifact and remains usable on real
rows. In the synthetic physical-lag experiment, the centered estimator recovers
\synRecover{} of the injected lag, consistent with the predicted attenuation of
\synLambda{} for this noise.

\section{The output-time specification sweeps}
\label{sup:sweeps}

\paragraph{Lag Recovery under the Full Control Set.} Adding
$-\tau_{\mathrm{inj}}\tilde v$ to the observed residual creates a known
perturbation along the fitted temporal regressor. The coefficient is recovered
at \injRecovery{} of its injected value at \injLevels\,ms, showing that the
geometric controls do not absorb this regressor. Injecting instead a purely
geometric displacement of \injGeoGamma{} px per px of box height moves the
minimal specification's $\tau$ by \injGeoMin\,ms and the reported one's by less
than \injGeoFull\,ms.

\paragraph{Sensitivity to Axis Pairing.} The reported specification
pairs box height and image row with $\hat{y}$ and box width and image column with
$\hat{x}$. Entering all four on both axes gives
$\satTau\,\pm\,\satTauSE$\,ms. Dropping each in turn moves $\tau$ by under
0.3\,ms except for the image column, whose removal returns the placebo to its
inflated value. Raising the overlap floor to \satTightIoU{} gives
$\satTight\,\pm\,\satTightSE$\,ms, the tightest of the sweep. Across every
specification carrying geometric controls, including speed floors from 5 to
20\,px/s, $\tau$ spans \satLo{} to \satHi\,ms, and none of them reaches the
occlusion-sensitivity centroid.

\paragraph{The bootstrap.} Resampling sequences with replacement,
\outBootDraws{} times, gives a \bootCI\,\% interval of
[\outBootLo,\,\outBootHi]\,ms with \outBootBeyond{} draws reaching the influence
centroid. The fit absorbs the per-sequence constant vectors by within-cell
demeaning, which reproduces the full design’s coefficient to \fwlAgree\,ms.

\section{The per-object ratio}
\label{sup:ratio}

For comparison with the estimator in main Eq.~2, we evaluate the per-object ratio
$-\langle\hat{\mathbf c}-\mathbf c^{*},\hat{\mathbf u}\rangle/\lVert \tilde{\mathbf v}\rVert$
on the same \outMatches{} rows. Across these rows, it has a sample standard deviation
of \ratioSD\,ms and an
interquartile range of \ratioIQR\,ms. Its mean over all rows is
$\ratioMean\,\pm\,\ratioMeanSE$\,ms, close to the $+\rvtCentroidMagMeas$\,ms coefficient
corresponding to the newest-window occlusion-sensitivity centroid under the convention
of main Eq.~2. Above a \ratioFloorSpeed\,px/s floor it gives
$\ratioFloor\,\pm\,\ratioFloorSE$\,ms, significant at more than two standard
errors, which is a fixed spatial bias divided by a speed: the ratio has no
intercept and main Eq.~2 does.

\section{The port of the released checkpoint}
\label{sup:port}

The released \texttt{rvt-t} checkpoint loads into a model built from the
repository's own configuration files with \rvtMissingKeys{} missing and
\rvtUnexpectedKeys{} unexpected keys at \rvtParams\,M parameters, and runs
forward under a current framework without the training code it was serialized
from. Both predictor measurements run on that port.

\begin{figure}[t]
\centering
\figasset{figs/fig3_day_night.pdf}{\columnwidth}
\caption{Left: within-frame dispersion of $\bar{t}_{i}$ against its analytic null
in the two exposure strata, \sigmaEvtFramesNight{} frames of
\texttt{zurich\_city\_09\_a} at \dsecExpMax\,\micro s and
\sigmaEvtFramesDay{} frames of \texttt{interlaken\_00\_c} at
\sigmaEvtExpDay\,\micro s. The daytime frames are selected by retaining frames
with at least three qualifying boxes. Right: per-pixel Rayleigh statistic $Z=R^{2}/N$ at
\flickerLineHz\,Hz and at \flickerCtrlHz\,Hz for the same two sequences, median
with a whisker to the ninety-fifth percentile, against the $\mathrm{Exp}(1)$
null. Both contrasts run in the same direction, and
Main Sec.~4.4 reports the controls applied to that dispersion.}
\label{fig:daynight}
\end{figure}


\begin{table}[t]
\centering\footnotesize
\setlength{\tabcolsep}{3pt}\renewcommand{\arraystretch}{0.95}
\begin{tabular}{lrrrr}
\toprule
Arm & frames & sd & null & excess \\
\midrule
Ceiling sequence, all pixels    & \sigmaEvtFramesNight & \sdAll & \nullAll & \sigmaEvtExcessNight \\
Phase-locked pixels removed     & \framesEleven & \sdClean & \nullClean & \excessClean \\
Random pixels, matched size     & \framesRandCtl & \sdRandCtl & \nullRandCtl & \excessRandCtl \\
Random-position boxes           & \framesRandBoxes & \sdRandBoxes & \nullRandBoxes & \excessRandBoxes \\
\midrule
\multicolumn{5}{l}{\emph{by each box's own \flickerLineHz\,Hz modulation depth $m$}}\\
$m\in[\modQoneLo,\modQoneHi)$    & \framesModQone & \sdModQone & \nullModQone & \excessModQone \\
$m\in[\modQoneHi,\modQtwoHi)$    & \framesModQtwo & \sdModQtwo & \nullModQtwo & \excessModQtwo \\
$m\in[\modQtwoHi,\modQthreeHi)$  & \framesModQthree & \sdModQthree & \nullModQthree & \excessModQthree \\
$m\in[\modQthreeHi,\modQfourHi]$ & \framesModQfour & \sdModQfour & \nullModQfour & \excessModQfour \\
\bottomrule
\end{tabular}
\caption{Per-object event centroid inside one published exposure window of
\texttt{zurich\_city\_09\_a}, in \micro s, each entry the median over frames of a
within-frame statistic against that frame's own null $w^{2}/12N_{i}$. The first
block's excess follows the main Sec.~4.4 statistic summarized over frames. Rows two and
three remove the \excessMaskPix{} pixels phase-locked at \flickerLineHz\,Hz and a
random subset of that size, row four places boxes of the same shapes at random
positions. In the lower block a frame enters a stratum when \minObjects{} of its
boxes fall in it, so those frame counts do not partition the first row's.}
\label{tab:controls}
\end{table}


\section{The DSEC measurements}
\label{sup:evidencetime}

\begin{figure}[t]
\centering
\figasset{figs/fig5_qualitative.pdf}{0.50\columnwidth}
\caption{Every event inside one published \dsecExpMax\,\micro s exposure window
of \texttt{zurich\_city\_09\_a}, colored by its position within that window, with
that frame's released boxes numbered left to right. The frame is the
median-dispersion one of \figFiveCands{} eligible candidate frames with at least
six qualifying boxes, not the largest.}
\label{fig:qualitative}
\end{figure}

\begin{table}[t]
\centering\footnotesize
\setlength{\tabcolsep}{3.6pt}\renewcommand{\arraystretch}{0.95}
\begin{tabular}{lrrrrr}
\toprule
Sequence & boxes & $m$ & $Z_{\flickerLineHz}$ & $Z_{\flickerCtrlHz}$ & locked \\
\midrule
\texttt{zurich\_city\_00\_a} & \sixBoxesA & \sixModA & \sixZA & \sixZcA & \sixLockA\,\% \\
\texttt{zurich\_city\_01\_a} & \sixBoxesB & \sixModB & \sixZB & \sixZcB & \sixLockB\,\% \\
\texttt{zurich\_city\_02\_a} & \sixBoxesC & \sixModC & \sixZC & \sixZcC & \sixLockC\,\% \\
\texttt{zurich\_city\_03\_a} & \sixBoxesD & \sixModD & \sixZD & \sixZcD & \sixLockD\,\% \\
\texttt{zurich\_city\_09\_a} & \sixBoxesE & \sixModE & \sixZE & \sixZcE & \sixLockE\,\% \\
\texttt{zurich\_city\_10\_a} & \sixBoxesF & \sixModF & \sixZF & \sixZcF & \sixLockF\,\% \\
\bottomrule
\end{tabular}
\caption{Every DSEC training sequence whose exposure is pinned at
\dsecExpMax\,\micro s, measured per labeled box over the events inside that box
during that frame's own published exposure window, with at least
\sixMinEvents{} events per box. $m$ is the median over boxes of the modulation
depth $m=2|C|$ at \flickerLineHz\,Hz. $Z_{\flickerLineHz}$ and
$Z_{\flickerCtrlHz}$ are the median
over boxes of the Rayleigh statistic at \flickerLineHz\,Hz and at the
off-frequency control on the same events. The analytic $\mathrm{Exp}(1)$ median of
\flickerZnull{} does not describe either column, which is why the two are compared
with each other. The locked fraction denotes the fraction of boxes whose local excess at
\flickerLineHz\,Hz exceeds the ninety-ninth percentile of the same quantity at
\flickerCtrlHz\,Hz.}
\label{tab:sixseq}
\end{table}

\paragraph{Data path.} For each labeled frame of a sequence, the frame's own
exposure interval $[a,b]$ is read from the dataset's published
\texttt{image\_exposure\_timestamps\_left} file, the event stream is sliced to
that interval through the release's \texttt{ms\_to\_idx} index, and the events
inside each released box are collected. The \texttt{ms\_to\_idx} index resolves to
whole milliseconds, so the slice it returns is filtered to $a\le t<b$ in
microseconds before any statistic is formed, and the analytic null is computed for
the same interval. Events and frames share one clock in this dataset, so no
calibration step enters. For the dispersion of main Sec.~4.4 a box contributes
only if it carries at least \minEvents{} events in the window, and a frame only
if at least \minObjects{} boxes qualify, the reported statistics are medians over
frames of a within-frame quantity, so no frame is weighted by its object count.
For the per-box statistics of Table~\ref{tab:sixseq} a box contributes with at
least \sixMinEvents{} events and no frame-level requirement applies. The selected
qualitative frame is shown in Supplement Fig.~\ref{fig:qualitative}, and the paired
daytime control is shown in Supplement Fig.~\ref{fig:daynight}.

For frame $f$, let $I_f$ be the qualifying boxes, let $N_{fi}$ be the event count
inside box $i$, and let $\bar t_{fi}=N_{fi}^{-1}\sum_{k\in i}t_k$ be its event-time
centroid relative to the frame's mid-exposure. The observed within-frame dispersion
and its frame-specific analytic null are
\[
\begin{aligned}
d_f&=\left[\frac{1}{|I_f|-1}\sum_{i\in I_f}
 (\bar t_{fi}-\bar t_f)^2\right]^{1/2},\\
n_f&=\left[\frac{1}{|I_f|}\sum_{i\in I_f}
 \frac{w_f^2}{12N_{fi}}\right]^{1/2}.
\end{aligned}
\]
where $\bar t_f=|I_f|^{-1}\sum_{i\in I_f}\bar t_{fi}$ and $w_f=b_f-a_f$.
The frame-level excess is $e_f=[\max(d_f^2-n_f^2,0)]^{1/2}$, and the reported
entries are the medians of $d_f$, $n_f$, and $e_f$ over eligible frames.

\paragraph{The stereo exposures.} The left and right exposures open at the same microsecond
in \dsecCotriggerPct\,\% of the \dsecCotriggerFrames{} frames checked and differ only in
when they close, and the published image timestamp equals the mean of the two
mid-exposures to within \dsecTsResidual\,\micro s.

\paragraph{Archive version and track association.} DSEC-Det ships more than
one annotation version. The numbers in main Sec.~4.1 describe the archive
\texttt{dsec-det\_left\_object\_detections.zip}, \ddArchiveBytes{} bytes,
retrieved \ddArchiveDate{}. Whether those labels are tracker output on the RGB
frames, which the dataset's own description states, is not distinguishable from
human annotation by a second-difference residual. On the Gen1 side the
\texttt{track\_id} field is a per-frame index rather than a persistent identifier,
so association for the output-time fit is rebuilt by same-class overlap between
adjacent label times, at which \outLinkPct\,\% of boxes link.

\paragraph{Null.} For events uniform on $[a,b]$ the mean of $N_{i}$ times has
variance $w^{2}/12N_{i}$ with $w=b-a$. Each frame therefore carries its own null,
formed from the published width and the observed counts, and the excess of main
Sec.~4.4 is summarized over frames rather than formed from the medians of the two
summaries. That is why the excess column of Table~\ref{tab:controls} is not
$(\mathrm{sd}^{2}-\mathrm{null}^{2})^{1/2}$ of the two columns beside it.

\paragraph{Phase test.} The per-pixel test at frequency $f$ takes the event phases
$\phi_{k}=2\pi f t_{k} \bmod 2\pi$ of a pixel, forms the resultant length $R$ and
reports $Z=R^{2}/N$, which is $\mathrm{Exp}(1)$ under uniform phase. Raw
modulation depth is not usable at these counts, because for $N$ events at random
phase the resultant is already of the size of the effect. Pixels enter with at
least \flickerMinEvents{} events over the tested slice, and every fourth event is
used to reduce sample size. The off-frequency and daytime controls remain close
to their nominal reference under the same procedure. Two arms hold the estimator to its null, an
off-frequency arm on the same events and the same frequency on daytime events.
The fraction of locked pixels is reported relative to live pixels in main Sec.~4.2 and
separately relative to the full sensor. The two denominators differ, and each statement
identifies its denominator.

\paragraph{The per-box statistics of Table~\ref{tab:sixseq}.} The same estimator is applied
to the events inside one released box during one published exposure window rather
than to the events of one pixel over a long slice. For box $i$ the depth is
$m=2|C|$ and the statistic is $Z=N_{i}|C|^{2}$ with
$\omega=2\pi f$ and $C=N_{i}^{-1}\sum_{k\in i}e^{-\mathrm{i}\omega t_{k}}$, phases referred to that
frame's mid-exposure instant. A finite exposure does not produce uniform phase
at a tested frequency, so the pixel-level $\mathrm{Exp}(1)$ null is not assumed
for these per-box values. The depth is reported descriptively, without a
random-phase floor. Boxes are matched to exposure windows by nearest mid-exposure
instant within half a frame period. The off-frequency arm at
\flickerCtrlHz\,Hz is run on exactly the same events.

\paragraph{Pixel exclusion and modulation strata.} The exclusion mask of
Table~\ref{tab:controls} is the set of pixels reaching $p<10^{-3}$ in the
per-pixel test, and the
matched control draws a random pixel set of exactly the same size, so the two
arms remove the same number of pixels. The mask is built from a six-second slice
at every fourth event. The stratification uses no mask: the quartiles of $m$ are
taken over all \modBoxes{} qualifying boxes, and a frame contributes to a stratum
when at least \minObjects{} of its boxes fall inside it, which is why the frame
counts of that block are not a partition of the first row's.

\section{The permutation null}
\label{sup:perm}

The analytic null $w^{2}/12N_{i}$ for the main Sec.~4.4 statistic is the variance of a mean of
$N_{i}$ draws uniform on the exposure, so it assumes the events of a box arrive
independently and uniformly.

The permutation null does not assume independent or uniform event arrivals. Instead,
it conditions on the pooled frame-level event times and each box's event count while
randomizing box assignment. For each frame the events of all qualifying boxes are pooled
and reassigned to boxes at random, preserving each box's count. The exposure, the flicker and the scene's burstiness
all survive that reassignment, only which box an event belongs to is destroyed.
Over \permFrames{} frames and \permDraws{} draws each it returns
\permNull\,\micro s, \permRatio{} times the analytic null rather than above it,
and \permAbovePct\,\% of frames exceed their own permutation null. A burst common to the frame
displaces every $\bar{t}_{i}$ in the same direction and cancels in their
dispersion, which is why the statistic was already insensitive to it. For a separate
random-position draw evaluated against the permutation null, boxes of the same shapes give
\permRandObs\,\micro s against \permRandNull\,\micro s. Thus the dispersion is not specific
to annotated object regions, although object regions may still affect it.

\section{The carried state of the SSM release}


Main Sec.~3.3 measures five released checkpoints of two architectures. The second is
SSM-ViT (Zubic et al., CVPR 2024), a fork of the RVT repository in which each stage's
ConvLSTM is
replaced by an S5 layer. It distributes RVT's own preprocessed Gen1 and builds its window
boundaries with the same statement of the same script, so the window convention of main
Sec.~3.2 holds there unchanged.

Its recurrent state does not. Let $s$ denote the carried state entering the first
step of a chunk. That release injects it as
$\bar{\Lambda}_{0}\leftarrow\bar{\Lambda}_{0}\odot s$ before an associative scan whose
binary operator is $(a_{i},b_{i})\circ(a_{j},b_{j})=(a_{j}a_{i},\,a_{j}b_{i}+b_{j})$. The
scan's outputs are the $b$ components, and $b_{p}$ is built from $a_{1},\dots,a_{p}$ and
$b_{0},\dots,b_{p}$, $\bar{\Lambda}_{0}$ enters only the $a$ component of the first element
and therefore no output. The carried state is inert.

We verify this code-level property empirically. Driving each model the way its own release
drives it and zeroing the state entering a chunk of \ssmChunkMeas{} windows, the relative
$L_{2}$ change in the emitted detection tensor is \ssmStateEffect{} at every position for
SSM-ViT, and \rvtStateEffectOne{} at position one falling to \rvtStateEffect{} at the last
position for \texttt{rvt-t}. A model that ignored its history altogether would return zero to
both arms, so the control occludes the window one position earlier inside the same chunk:
SSM-ViT answers \ssmInchunkOne{} there, falling to \ssmInchunkLast{}, so it is recurrent
within a chunk and only across chunks is the state lost.

\section{Effect of Temporal Displacement on Checkpoint Ordering}

Main Sec.~3.5 reports that displacing the ground truth to the newest window's occlusion-sensitivity centroid costs
\mapCost{} points on one checkpoint. Whether the interval changes a \emph{comparison} needs
more than one detector, so all \rankCkpts{} released checkpoints of main Sec.~3.3 were run
over the Gen1 validation split under one protocol: the same frames, the same confidence
threshold and non-maximum suppression, and ground truth that is identical across the five and
identical to the dump main Sec.~3.5 already uses. Each was then scored at the label instant
and again at the displacement it prefers, which ranges from \rankArgLo{} to \rankArgHi\,ms
across the five.

\begin{table}[h]
\centering\footnotesize
\setlength{\tabcolsep}{3.4pt}\renewcommand{\arraystretch}{0.95}
\begin{tabular}{ll}
\toprule
stratum & ordering at the label instant \\
\midrule
all moving       & b $>$ s $>$ t $>$ S5-B $>$ S5-S \\
$10$--$25$\,px/s & S5-B $>$ b $>$ t $>$ S5-S $>$ s \\
$25$--$50$\,px/s & S5-B $>$ b $>$ S5-S $>$ s $>$ t \\
$\geq 50$\,px/s      & b $>$ t $>$ S5-S $>$ s $>$ S5-B \\
\bottomrule
\end{tabular}
\caption{Ordering of the \rankCkpts{} released checkpoints at the label instant, by object
speed, \texttt{t}, \texttt{s}, \texttt{b} are the RVT capacities and S5-S, S5-B the SSM ones. Scoring each at its own preferred displacement leaves every one of these
orderings unchanged. A common displacement sweep reveals one localized inversion
between \texttt{rvt-b} and S5-S in the 25--50\,px/s stratum. Ranking varies across
speed strata, and the common sweep produces no other pairwise inversion. \rankMover{}
is first in the middle two strata and last in the fastest, and \texttt{rvt-t} is third
overall, last between \mapCutMid{} and \mapCutHi\,px/s and second above \mapCutHi{}.}
\end{table}

The own-preferred-offset scores preserve the reported ordering in each speed
stratum. Under a common displacement sweep, the $25$--$50$\,px/s stratum
contains a localized ranking inversion between \texttt{rvt-b} and S5-S. The
other evaluated strata contain no sign-changing pair. The aggregate ranking and
the fast-object ranking have different leaders, so the strata are reported
separately.

The released streaming evaluation steps through the split in non-overlapping chunks of
\ssmEvalChunk{} windows (\texttt{sequence\_length} in \texttt{config/dataset/gen1.yaml},
\texttt{start\_indices} in \texttt{sequence\_for\_streaming.py} steps by exactly that). The
evaluated input sensitivity of a released detection therefore depends on where its frame falls inside
that chunk, from a full \ssmEvalChunk{} windows at the last position to the current window
alone at the first. This is reported as a property of the release that its documentation does
not state, which is the same class of undeclared quantity the rest of this paper measures,
and not as a claim about what the architecture could support if the state were carried.

\end{document}
